% Analytical checkpoint: post-foundations progress toward Theorems A and B.
% Version 1.0. Sources are listed in the provenance appendix.
% Do not append new research directly to this file. Deliver review increments first.
\documentclass[11pt]{article}
\usepackage[T1]{fontenc}
\usepackage[utf8]{inputenc}
\usepackage{lmodern}
\usepackage[margin=0.9in]{geometry}
\usepackage{amsmath,amssymb,amsthm,mathtools,mathrsfs}
\usepackage{booktabs,longtable,array,enumitem,microtype}
\usepackage{xurl}
\usepackage[hidelinks]{hyperref}
\usepackage{fancyhdr}
\pagestyle{fancy}
\fancyhf{}
\fancyhead[L]{\small Analytical progress toward Theorems A and B}
\fancyhead[R]{\small Review checkpoint 1.0}
\fancyfoot[C]{\thepage}
\setlength{\headheight}{14pt}
\setlength{\emergencystretch}{3em}
\allowdisplaybreaks[2]
\numberwithin{equation}{section}
\newtheorem{proposition}{Proposition}[section]
\newtheorem{lemma}[proposition]{Lemma}
\newtheorem{theorem}[proposition]{Theorem}
\newtheorem{corollary}[proposition]{Corollary}
\theoremstyle{definition}
\newtheorem{definition}[proposition]{Definition}
\newtheorem{target}{Open target}[section]
\newtheorem{calculation}[proposition]{Calculation}
\theoremstyle{remark}
\newtheorem{remark}[proposition]{Remark}
\newcommand{\E}{\mathbb E}
\newcommand{\R}{\mathbb R}
\newcommand{\1}{\mathbf 1}
\newcommand{\dd}{\,d}
\newcommand{\norm}[1]{\left\lVert #1\right\rVert}
\newcommand{\abs}[1]{\left|#1\right|}
\DeclareMathOperator{\Cov}{Cov}
\DeclareMathOperator{\Var}{Var}
\DeclareMathOperator{\tr}{tr}
\DeclareMathOperator{\diag}{diag}
\newcommand{\scope}[1]{\par\noindent\textbf{Scope and status:} #1\par\smallskip}
\newcommand{\openobligation}[1]{\par\noindent\textbf{Open obligation:} #1\par\smallskip}
\newcommand{\source}[1]{\par\noindent\textit{Source:} #1\par\smallskip}
\title{Analytical progress toward initialized population reversal\\
\large Post-foundations checkpoint for Theorems A and B}
\author{Research working document}
\date{Version 1.0 --- consolidation for review}
\begin{document}
\maketitle

\begin{abstract}
This document consolidates the analytical progress made after
\path{RELU_POPULATION_FOUNDATIONS_v2.tex}. It records exact finite-noise
identities, initialized early-escape estimates, the noise-scaled population
system, an analytical reversal theorem for its distinguished two-atom orbit,
population compensation-lag and splitting results, and the exact-lag and
inactive-bridge refinements. \textbf{The initialized Theorems A and B for the
original Gaussian population are not yet proved.} Their entry, nonlinear
passage, large-bridge comparison, signed finite-noise control, and parameter
coverage remain explicit open obligations. The proof program is analytical:
no saved numerical trajectory, interval arithmetic, or numerical sign check
is a premise of a theorem here. Numerical observations appear only in a
separate evidence ledger.
\end{abstract}

\section*{Read-first contract}
\addcontentsline{toc}{section}{Read-first contract}
\textbf{This is a review checkpoint, not a submission-ready completed theory.}
A proof below establishes only its displayed statement under its displayed
hypotheses. A theorem within the \emph{defined leading system} is not silently
promoted to a theorem about the exact finite-noise population.
The source-derived formal reduction and the response calculations are
identified separately from proved identities and conditional implications.

The foundations, the original research notebook, and the old appendix remain
unchanged. New work must be delivered in separate, versioned proof-increment
files and checked before the user authorizes incorporation into a new version
of this checkpoint. Do not automatically \verb|\input| unreviewed increments.
Do not access or modify GitHub; the user handles version control.

The analytical parts of \cite{E}, Sections 1--3, are retained. Its numerical
reference construction and certificate roadmap are \emph{not} part of the
current proof program. The exact finite-parameter identities of \cite{F,LC}
remain useful, but their earlier suggestions to certify saved trajectories
are superseded by the analytical-only handoff.

\tableofcontents
\clearpage

\section{Sources, theorem objectives, and status ledger}\label{pa:status}
\source{The foundations \cite{F}; early escape \cite{E}; factorization audit
\cite{LC}; limiting orbit \cite{TN}; scaled population \cite{SP}; the subsequent
finite-noise-lag and bridge-feedback discussion \cite{C}.}

Theorem A is the \emph{initialized occurrence and exact attribution} target:
substantial learning of both concepts, a common learned interval with
$D_1<0<D_2$, and early-to-late dominance reversal uniformly over an off-cone
sector and a stated parameter region. Theorem B is the \emph{mechanism}
target: reciprocal leakage compensation, loss of probe exposure through
input rotation, and the effect of the probe-inactive bridge account for
those signs and their changing balance. The same analytical population
argument should supply both results.

\begin{longtable}{@{}p{0.28\textwidth}p{0.30\textwidth}p{0.35\textwidth}@{}}
\toprule
Component & Current mathematical status & Not implied by this component\\
\midrule
\endhead
Finite-noise foundations & Imported from \cite{F}; exact flow, rates,
regularity, balance, and linear control & Initialized learned-state entry or
canonical late signs\\
Early escape, Section~\ref{pa:escape} & Analytical estimates from the actual
isotropic initialization & Substantial weak learning or late crossover\\
Exact residual bookkeeping, Section~\ref{pa:bookkeeping} & Identities and
conditional tracking estimates & Signs of unbounded forcing terms\\
Scaled population, Section~\ref{pa:scaled} & Leading equations derived;
internal results proved for that defined system & A uniform finite-$q$
trajectory approximation theorem\\
Resident and shape analysis, Section~\ref{pa:shape} & Exact resident curve,
linearized spectra and multipliers; same-field ordering & Nonlinear saddle
passage, isotropic capture, or a two-dimensional population unstable manifold\\
Population lag, Section~\ref{pa:lagpop} & Positive-lag law and helpful sign
within the leading $M=1$ system & Exact finite-$q$ lag positivity\\
Selected-orbit reversal, Section~\ref{pa:orbit} & Analytical reversal on the
coherent weak-growth branch & Selection of that branch by isotropic data,
canonical constants, or a specified weak-learning threshold\\
Exact lag and bridge, Section~\ref{pa:finiteq} & Exact finite-$q$ state
identities and leading-system instantaneous bridge signs & Between-population
trajectory ordering or canonical signed-rate closure\\
Small-split response, Section~\ref{pa:response} & Response functional and
second-order calculation & An analytical proof of $\kappa>0$ or extension to
the canonical inactive bridge\\
Initialized A and B & \textbf{Open}; roadmap in Section~\ref{pa:roadmap} & No
completion claim is made\\
\bottomrule
\end{longtable}

\paragraph{Consolidation choices.}
All derivatives use normalized time $\tau$ unless explicitly marked physical.
The uniform label law is $\lambda_0$, whereas $\lambda=\rho^2$ is a growth
rate. The weak resident coordinate is allowed to be real; this incorporates
the already-discussed removal of the artificial $\rho\le1/\sqrt2$ restriction.
The mean lag, the $X_1$-weighted lag, and the leading scaled lag receive
different symbols. The early-escape source's temporary $u/\varepsilon_{
\rm al}$ naming inconsistency is resolved by defining $\varepsilon_{\rm al}$
explicitly. These are disclosed notation/scope reconciliations, not new
initialized proof claims.

\section{Interface imported from the foundations}\label{pa:setup}
For $0<\rho<1$ and $q>0$, take
\begin{equation}\label{pa:data}
 P_1=\mathcal N(e_1,q^2I_2),\qquad
 P_2=\mathcal N(\rho e_2,q^2I_2),\qquad P=(P_1+P_2)/2.
\end{equation}
The task is identity regression, with no origin cluster and no biases.
Both layers of the ReLU network are trained independently. Put
\[
 a_\theta=(\cos\theta,\sin\theta),\quad
 b_\psi=(\cos\psi,\sin\psi),\quad
 \lambda_0(d\alpha)=\frac{d\alpha}{2\pi}.
\]
The balanced characteristic representation is
\begin{equation}\label{pa:representation}
 U_\alpha=\sqrt{m_\alpha}a_{\theta_\alpha},\quad
 W_\alpha=\sqrt{m_\alpha}b_{\psi_\alpha},\quad
 \nu=\int m_\alpha\delta_{(\theta_\alpha,\psi_\alpha)}\dd\lambda_0,
 \qquad f(X)=\int b(a^\top X)_+\dd\nu.
\end{equation}
Initialization is exactly
\begin{equation}\label{pa:init}
 \theta_0(\alpha)=\psi_0(\alpha)=\alpha,\qquad m_0(\alpha)=S_0.
\end{equation}
In finite-width coordinates this corresponds to
$u_i(0)=w_i(0)=\sqrt{S_0/h}\,a_{\alpha_i}$ and
$S_0=h\varepsilon_h^2$. Balance means equality of norms; vector alignment
need not persist. In the continuum,
\[
 f_0(X)=\frac{S_0}{4}X,\qquad E_\xi(0)=\frac12(1-S_0/4)^2
 \quad (|\xi|=1).
\]
With $\rho=\mu_2/\mu_1$, $q=\sigma/\mu_1$, use
$\tau=\mu_1^2t_{\rm phys}$. There is no rescaling of $S_0$ and no additional
width factor in this continuum time. For a fixed numerical unit probe,
$D_p^{\rm phys}=\mu_1^2D_p^\tau$. Rescaling the probe itself by $\mu_1$
would introduce an additional squared-amplitude factor and is not done here.

Set $e_f(X)=X-f(X)$ and
\begin{equation}\label{pa:residualmatrix}
 R_{p,f}(\theta)=\E_{P_p}[e_f(X)X^\top\1_{\{a_\theta^\top X>0\}}],
 \qquad R_f=(R_{1,f}+R_{2,f})/2.
\end{equation}
The imported exact equations are
\begin{equation}\label{pa:exactflow}
 \theta'=b^\top R_fa^\perp,\qquad
 \psi'=(b^\perp)^\top R_fa,\qquad
 (\log m)'=2b^\top R_fa.
\end{equation}
For a unit probe $\xi$, let $e_\xi=\xi-f(\xi)$ and
\begin{align}
 V_p(\xi)&=\int\left[(a^\top\xi)_+R_pa+
 \1_{\{a^\top\xi>0\}}bb^\top R_p\xi\right]\dd\nu,\label{pa:rawvelocity}\\
 D_p^\tau(\xi)&=-\frac12e_\xi^\top V_p(\xi),\qquad
 E_\xi'=D_1^\tau+D_2^\tau.\label{pa:exactrates}
\end{align}
The derivative identity is interpreted almost everywhere in time, with the
probe-boundary convention and regularity specified in \cite{F}. For pointwise
strict rate continuity, boundary mass must be excluded or explicitly controlled.
The training-cluster index does not name a neuron cohort.

Partition-free learning observables are
\begin{equation}\label{pa:learning}
 \mathsf m_p=\frac{\E_{P_p}[X_pf_p(X)]}{\E_{P_p}X_p^2},\qquad
 \mathcal L_p=\frac12\E_{P_p}|f(X)-X|^2.
\end{equation}
Neither family mass nor decreasing total loss alone proves the required
learning and signed-residual conclusions. The exact linear control
$D_p^{\rm lin}\le0$ is already Proposition P10 of \cite{F}; it is not a new
result of this checkpoint.

\section{Analytical initialized early escape}\label{pa:escape}
\source{\cite{E}, Sections 1--3 only. The later numerical-reference construction
is excluded from the present program.}

\subsection{Spectral and alignment control}
Put
\[
 A=\E_P[XX^\top]=\diag(1/2+q^2,\rho^2/2+q^2),\qquad a_*=1/2+q^2,
\]
\[
 A_\theta=\E_P[XX^\top\1_{\{a^\top X>0\}}],\qquad
 C_f(\theta)=\E_P[f(X)X^\top\1_{\{a^\top X>0\}}].
\]
Thus $R_f=A_\theta-C_f$. Define
\begin{align}
 \overline S(\tau)&=S_0e^{2a_*\tau},&
 \varepsilon_{\rm al}(\tau)&=\overline S(\tau)-S_0,\label{pa:earlyenvelopes}\\
 I(\tau)&=a_*\int_0^\tau\varepsilon_{\rm al}(s)^2\dd s
 =S_0^2\left[\frac{e^{4a_*\tau}-1}{4}
 -(e^{2a_*\tau}-1)+a_*\tau\right].\nonumber
\end{align}

\begin{proposition}[Initialized spectral, mass, and alignment bounds]
\label{pa:earlybounds}
For the exact initialized flow, $S(\tau)=\int m_\alpha\dd\lambda_0\le
\overline S(\tau)$. On every interval where
$\varepsilon_{\rm al}(\tau)<\pi$, continuous angular lifts satisfy, uniformly
in the initial label,
\begin{equation}\label{pa:earlybounds-eq}
 |\psi-\theta|\le\varepsilon_{\rm al},\qquad
 S_0e^{-\varepsilon_{\rm al}-I}\le m_\alpha
 \le S_0e^{2a_*\tau+\varepsilon_{\rm al}}.
\end{equation}
\end{proposition}
\scope{Exact finite-noise analytical statement, initialized at \eqref{pa:init}.}
\begin{proof}
Positivity gives $0\preceq A_\theta\preceq A\preceq a_*I$. For unit vectors
$v,z$, the feature representation and Cauchy--Schwarz give
\[
 |v^\top C_fz|
 \le\int m_\alpha\E[(a_{\theta_\alpha}^\top X)_+|z^\top X|]
 \dd\lambda_0\le a_*S.
\]
The exact mass identity yields
\[
 S'=2\E[X^\top f]-2\E|f|^2\le2a_*S,
\]
because $\E[X^\top f]=\int m b^\top A_\theta a\dd\lambda_0\le a_*S$.
Integrating proves the envelope.
For $\delta=\psi-\theta$, every symmetric planar matrix $B$ satisfies
\[
 (b^\perp)^\top Ba-b^\top Ba^\perp=-\tr(B)\sin\delta.
\]
Using $R=A_\theta-C_f$ therefore gives
\[
 \delta'=-\tr(A_\theta)\sin\delta+\epsilon_\delta,
 \qquad |\epsilon_\delta|\le2a_*S.
\]
Before a possible first exit from $|\delta|<\pi$, the target term points
inward and $D^+|\delta|\le2a_*\overline S$. Integration gives
$|\delta|\le\varepsilon_{\rm al}$, and closes the first-exit condition.
Finally,
\[
 b^\top A_\theta a
 =\tfrac12\{a^\top A_\theta a+b^\top A_\theta b
 -(b-a)^\top A_\theta(b-a)\}\ge-\tfrac12a_*\delta^2.
\]
Thus $-a_*\delta^2-2a_*S\le(\log m)'\le2a_*+2a_*S$.
Integrating and substituting the envelopes proves \eqref{pa:earlybounds-eq}.
\end{proof}

\subsection{Comparison with an exact target-only flow}
For each initial label define the mathematical comparison solution
\[
 \widehat V'=\E_P[X(\widehat V^\top X)_+],\qquad
 \widehat V_0=\sqrt{S_0}a_\alpha.
\]
Set $\widehat U=\widehat W=\widehat V$,
$\widehat m=|\widehat V|^2$, and let $\widehat\theta$ be its angle.
This is not a numerically supplied reference and is not asserted to equal
the nonlinear flow. With
\[
 v_0(\theta)=(a^\perp)^\top A_\theta a,\qquad
 g_0(\theta)=a^\top A_\theta a,
\]
its equations are $\widehat\theta'=v_0(\widehat\theta)$ and
$(\log\widehat m)'=2g_0(\widehat\theta)$.

\begin{proposition}[Target-only comparison]\label{pa:targetcomparison}
On the interval of Proposition~\ref{pa:earlybounds}, define
\[
 H(\tau)=S_0(2e^{2a_*\tau}-3e^{a_*\tau}+1).
\]
Then, labelwise,
\[
 |\theta-\widehat\theta|\le H,\qquad
 |\psi-\widehat\theta|\le H+\varepsilon_{\rm al},\qquad
 \left|\log\frac m{\widehat m}\right|\le2H.
\]
\end{proposition}
\begin{proof}
The moving-boundary identity gives $A_\theta'a=0$. Consequently
\[
 v_0'=(a^\perp)^\top A_\theta a^\perp-a^\top A_\theta a,
 \qquad g_0'=2(a^\perp)^\top A_\theta a,
\]
and $|v_0'|,|g_0'|\le a_*$. For the second bound, compare the quadratic form
on $(a+a^\perp)/\sqrt2$ and $(a-a^\perp)/\sqrt2$; both values lie in
$[0,a_*]$. The true velocity satisfies
\[
 |\theta'-v_0(\theta)|\le a_*(\varepsilon_{\rm al}+\overline S).
\]
The upper Dini derivative of $h=|\theta-\widehat\theta|$ is at most
$a_*h+a_*(\varepsilon_{\rm al}+\overline S)$. The displayed $H$ solves that
scalar equality with zero initial value. Moreover,
\[
 |b^\top R_fa-g_0(\widehat\theta)|
 \le a_*(\varepsilon_{\rm al}+\overline S+H)=H'.
\]
Integrate the logarithmic mass equations and then use the alignment bound.
\end{proof}

\subsection{A nontrivial canonical mass escape}
\begin{proposition}[Canonical early escape]\label{pa:canonicalescape}
For $\rho=2/3$, $q=1/20$, $S_0=1/5000$, on $0\le\tau\le6$,
\[
 S\le0.0832,\quad |\psi-\theta|\le0.0830,\quad
 1.83\cdot10^{-4}<m_\alpha<0.091,
 \quad |U_\alpha|=|W_\alpha|>0.0135.
\]
Furthermore $S(6)>16S_0=0.0032$.
\end{proposition}
\begin{proof}
Set $\theta_0=\pi/12$ and select the fixed initial arc
$C_0=[-\theta_0,\theta_0]$, of label probability $1/12$. Gaussian integration
expresses the target-only angular field as
\begin{align*}
 v_0(\theta)=\tfrac12\big[&-\sin\theta\{\cos\theta\Phi(\cos\theta/q)
 +q\phi(\cos\theta/q)\}\\
 &+\rho\cos\theta\{\rho\sin\theta\Phi(\rho\sin\theta/q)
 +q\phi(\rho\sin\theta/q)\}\big].
\end{align*}
At $\theta_0$, use $\Phi(\cos\theta_0/q)\ge1/2$, $\Phi\le1$ and
$\phi\le2/5$ to obtain
\[
 v_0(\theta_0)\le\frac{\cos\theta_0}{2}
 \left[-\frac{\sin\theta_0}{18}+\frac1{75}\right]<0.
\]
The last inequality uses $\sin(\pi/12)>1/4$. At $-\theta_0$, both terms in
the original expression are positive, so the comparison arc is invariant.
On it Jensen's inequality gives
\[
 g_0(\widehat\theta)=\E_P(a_{\widehat\theta}^\top X)_+^2
 \ge\tfrac12\cos^2\theta_0.
\]
Proposition~\ref{pa:targetcomparison} yields
\begin{equation}\label{pa:escapelower}
 S(\tau)\ge\frac{S_0}{12}
 \exp\{\cos^2(\pi/12)\tau-2H(\tau)\}.
\end{equation}
Here $a_*=201/400$. The elementary inequalities
$e^{6.03}<416$ and $e^{3.015}>20$ imply
$\overline S(6)<0.0832$, $\varepsilon_{\rm al}(6)<0.0830$, and
$H(6)<0.1546$. Also
\[
 I(6)\le S_0^2(416^2/4+3.015)=0.0017306806.
\]
The radius bounds follow from $e^{-r}\ge1-r$ and
$e^r\le(1-r)^{-1}$ for $0\le r<1$. Finally,
\[
 6\cos^2(\pi/12)-2H(6)>6(0.933)-0.3092=5.2888,
 \qquad e^{5.2888}>192.
\]
Equation~\eqref{pa:escapelower} proves the claim. These exponential
inequalities may be checked directly by finite rational Taylor sums: for
$x\ge0$ and $k+2>x$,
\[
 \sum_{j=0}^k\frac{x^j}{j!}
 \le e^x\le
 \sum_{j=0}^k\frac{x^j}{j!}
 +\frac{x^{k+1}}{(k+1)!}\frac1{1-x/(k+2)}.
\]
For example use $k=30$ for the first upper bound and sufficiently many
positive terms for the two lower bounds. This is elementary arithmetic in
the analytical proof, not a trajectory certificate.
\end{proof}
\openobligation{This stops at physical time $2/3$ when $\mu_1=3$.
It proves mass escape, not strong/weak chart entry, substantial weak learning,
late signed competition, or a parameter region for Theorem A.}

\section{Exact finite-parameter compensation and rate bookkeeping}
\label{pa:bookkeeping}
\source{\cite{LC} and the fixed-label compensation-tracking derivation in
\cite{C}. The equations below use normalized rates; multiply them by
$\mu_1^2$ for the physical rates in the diagnostic CSVs.}

\subsection{The full motion decomposition}
For a fixed probe set $z=a^\top\xi$ and $z_\perp=(a^\perp)^\top\xi$.
Define the source-specific velocities
\[
 g_p=\tfrac12b^\top R_pa,\qquad
 v_{\theta,p}=\tfrac12b^\top R_pa^\perp,\qquad
 v_{\psi,p}=\tfrac12(b^\perp)^\top R_pa.
\]
\begin{proposition}[Source-by-motion decomposition]\label{pa:motion}
At times for which the probe chain rule holds,
\begin{align}\label{pa:motion-eq}
 D_p^\tau=-\int\big[&2g_pz_+(e_\xi^\top b)
 +v_{\psi,p}z_+(e_\xi^\top b^\perp)\nonumber\\
 &+v_{\theta,p}\1_{\{z>0\}}z_\perp(e_\xi^\top b)\big]\dd\nu.
\end{align}
\end{proposition}
\begin{proof}
Differentiate $m b(a^\top\xi)_+$ along each labeled characteristic,
using $m'_p=2mg_p$, $b'_p=v_{\psi,p}b^\perp$ and
$a'_p=v_{\theta,p}a^\perp$. Contract the integrated output velocity
with $-e_\xi$. This also recovers \eqref{pa:exactrates}.
\end{proof}
No neuron family has been dropped. A moving cohort may be used to evaluate
this identity at a fixed time. Differentiating that cohort's mass or training
output later requires its membership flux, or a fixed-label construction.

\subsection{Gaussian projection compensation}
Let $A=\{a^\top\xi>0\}$ and $I=A^c$ at the current state, and write
$f=f^A+f^I$. These $A,I$ are population subsets, not the data moment matrix.
Put $Q_1=1+q^2$ and
\[
 c_A=\frac{\E_{P_1}[X_1f_2^A(X)]}{Q_1},\qquad
 \ell_I=\frac{\E_{P_1}[X_1f_2^I(X)]}{Q_1},\qquad
 r_{21}=\frac{\E_{P_1}[X_1(f_2(X)-X_2)]}{Q_1}.
\]
\begin{proposition}[Exact compensation projection]\label{pa:projection}
One has
\begin{equation}\label{pa:projection-eq}
 c_A=-\ell_I+r_{21},\qquad
 |r_{21}|\le
 \left(\frac{\E_{P_1}(f_2-X_2)^2}{Q_1}\right)^{1/2}.
\end{equation}
In particular, $\ell_I>|r_{21}|$ forces $c_A<0$.
\end{proposition}
\begin{proof}
The first identity follows from $f=f^A+f^I$ and
$\E_{P_1}[X_1X_2]=0$. The bound is Cauchy--Schwarz.
\end{proof}
\begin{remark}
This is a Gaussian-distribution projection, not merely the centroid value
$f_2(e_1)$. Neither the identity nor a small total loss proves a positive
signed lower bound on $r_{21}$.
\end{remark}

There is a useful exact transfer identity. Define
\[
 h_1(\theta)=\frac{\E_{P_1}[X_1(a_\theta^\top X)_+]}{Q_1}.
\]
For any scalar $k>0$, since $f^I(\xi)=0$,
\begin{equation}\label{pa:transferprojection}
 f_2(\xi)=-k\ell_I+kr_{21}
 +\int_A b_2[a^\top\xi-kh_1(\theta)]\dd\nu.
\end{equation}
The last term has absolute value at most
$\int_A|b_2|\,|a^\top\xi-kh_1(\theta)|\dd\nu$.
This proves how a compensating training projection can be exposed at a probe,
provided the displayed geometry error is bounded; it does not by itself
prove a sign of $E_\xi'$.

\subsection{Exact help and harm factorizations}
Choose a measurable core $C\subseteq A$ of mass $M_C>0$. A bar and
$\langle\cdot\rangle_C$ denote normalized mass averages on $C$.
For this subsection set
\[
 v=-z_\perp,\quad d=-e_\xi^\top b^\perp,\quad c=e_\xi^\top b,
 \quad G_1=R_1a,\quad T_2=R_2a^\perp,\quad
 \Lambda_1=-(b^\perp)^\top G_1,\quad j=b_1T_{2,1}.
\]
The quantities are signed, regardless of the absolute-value column names
used in diagnostic output. Define
\begin{align}
 B&=\frac{M_C}{2}\langle zd\Lambda_1\rangle_C,&
 H^{(1)}&=\frac{M_C}{2}\langle vcj\rangle_C,\label{pa:BH}\\
 H^{(2)}&=\frac{M_C}{2}\langle vc\,b_2T_{2,2}\rangle_C,&
 D_1^\tau&=-B+r_1,\quad
 D_2^\tau=H^{(1)}+H^{(2)}+r_2.\nonumber
\end{align}
The last equations define remainders containing all other motions and all
labels outside $C$. No smallness is assumed.

\begin{proposition}[Factorizations with complete errors]\label{pa:factorizations}
Let
\[
 B_0=\frac{M_C}{2}Q_1\bar z\,\bar d\,r_{21},\qquad
 H_{1,0}=\frac{M_C}{2}\bar v\,\bar c\,\bar j.
\]
Then $B=B_0+\epsilon_B$, $H^{(1)}=H_{1,0}+\epsilon_H$, where
\begin{align}
 \frac{2\epsilon_B}{M_C}
 &=Q_1r_{21}\Cov_C(z,d)
 +\langle zd(\Lambda_1-Q_1r_{21})\rangle_C,\label{pa:epsilonB}\\
 \frac{2\epsilon_H}{M_C}
 &=\Cov_C(vc,j)+\bar j\Cov_C(v,c).\label{pa:epsilonH}
\end{align}
In particular these are controlled by Cauchy--Schwarz and absolute weighted
moments, without dividing by a possibly vanishing mean residual.
\end{proposition}
\begin{proof}
For the first equation add and subtract
$Q_1r_{21}\langle zd\rangle_C$ in $\langle zd\Lambda_1\rangle_C$.
For the second use
$\langle vcj\rangle_C=\langle vc\rangle_C\bar j+\Cov_C(vc,j)$.
\end{proof}
The help reduction error is \emph{not only} within-core covariance. Put
\[
 R^{\rm full}_{1,22}=\E_{P_1}[e_2(X)X_2],\qquad
 \tau_1(\theta)=\E_{P_1}[e_2(X)(-a^\top X)_+].
\]
The identity $(a^\top X)_+=a_1X_1+a_2X_2+(-a^\top X)_+$ gives
\begin{align}\label{pa:helpdefect}
 \Lambda_1-Q_1r_{21}
 ={}&Q_1r_{21}(a_1b_1-1)-b_1a_2R^{\rm full}_{1,22}\nonumber\\
 &-b_1\tau_1(\theta)+b_2G_{1,1}(\theta).
\end{align}
This keeps mean tilt, transverse residual, wrong-gate tail, and first-coordinate
residual separately. For $a_1>0$,
\[
 |\tau_1(\theta)|\le\norm{e_2}_{L^2(P_1)}
 \big[(a_1^2+q^2)\Phi(-a_1/q)-a_1q\phi(a_1/q)\big]^{1/2}.
\]
The bound follows by Cauchy--Schwarz and the negative-part second moment of
a Gaussian with mean $a_1$ and variance $q^2$.

Define
\begin{equation}\label{pa:GLC}
 \mathfrak G_{\rm LC}=\bar v\,\bar c\,\langle b_1T_{2,1}\rangle_C
 -Q_1\bar z\,\bar d\,r_{21}.
\end{equation}
Then the exact bookkeeping identity is
\begin{equation}\label{pa:GLCidentity}
 E_\xi'=\frac{M_C}{2}\mathfrak G_{\rm LC}+\mathcal R,\qquad
 \mathcal R=\epsilon_H-\epsilon_B+H^{(2)}+r_1+r_2.
\end{equation}
A valid absolute budget is
$\mathcal E=|\epsilon_H|+|\epsilon_B|+|H^{(2)}|+|r_1|+|r_2|$.
Separate cluster signs require
\[
 B_0>|\epsilon_B|+|r_1|,\qquad
 H_{1,0}>|\epsilon_H|+|H^{(2)}|+|r_2|.
\]
This is not an autonomous reduced ODE, and no equality with the old archived
normal-form statistic $G$ has been proved.

\subsection{Source, carrier, and residual producer}
For a partition into output-producing groups $g$, write $f=\sum_g f^g$ and
\begin{align*}
 T_{2,1}^{\rm tar}(\theta)&=\E_{P_2}[X_1(a^\perp{}^\top X)
 \1_{\{a^\top X>0\}}],\\
 T_{2,1}^{g}(\theta)&=-\E_{P_2}[f_1^g(X)(a^\perp{}^\top X)
 \1_{\{a^\top X>0\}}].
\end{align*}
Linearity gives the exact identities
\begin{equation}\label{pa:producer}
 T_{2,1}=T_{2,1}^{\rm tar}+\sum_gT_{2,1}^g,\qquad
 H^{(1),g}=\frac12\int_C vc\,b_1T_{2,1}^g(\theta)\dd\nu.
\end{equation}
Cluster 2 can supply the gradient, the weak family's output can shape its
residual, and probe-active strong labels can carry the resulting motion.
These are three different indices. Evaluating $T_{2,1}^g$ at the mean angle
is not the rate-weighted integral in \eqref{pa:producer}.

\subsection{A conditional analytical compensation-tracking equation}
\begin{proposition}[Fixed-label tracking identity]\label{pa:tracking}
Let $A$ now be a \emph{fixed subset of initial labels}, $I=A^c$, and put
$h(X)=X_1/\sqrt{Q_1}$. Use the $L^2(P_1)$ inner product. Define
\[
 c=\langle h,f_2^A\rangle_1,\quad
 \ell=\langle h,f_2^I\rangle_1,\quad
 \delta=c+\ell=\langle h,f_2-X_2\rangle_1=\sqrt{Q_1}r_{21}.
\]
Let the output-weight kernel operator be
\[
 (\mathcal T_A r)(X)=\E_{P_1,X'}\!\left[
 \int_A(U_\alpha^\top X)_+(U_\alpha^\top X')_+\dd\lambda_0(\alpha)
 \,r(X')\right].
\]
Write the exact motion as
$\partial_\tau f_2^A=-\tfrac12\mathcal T_A(f_2-X_2)+v_A$, where $v_A$
contains all input motion of $A$ and its cluster-2 output-weight motion.
Define
\[
 \lambda_A=\langle h,\mathcal T_Ah\rangle_1,\quad
 k_A=\mathcal T_Ah-\lambda_Ah,\quad
 f_2-X_2=\delta h+r_\perp.
\]
Then
\begin{align}\label{pa:tracking-eq}
 \delta'&=-\kappa_A\delta+\eta_A,\qquad \kappa_A=\lambda_A/2\ge0,\\
 \eta_A&=\ell'+\langle h,v_A\rangle_1
 -\tfrac12\langle k_A,r_\perp\rangle_1.\nonumber
\end{align}
\end{proposition}
\begin{proof}
The kernel is self-adjoint positive semidefinite by its feature-integral
representation. Cluster-1 output-weight descent is exactly the displayed
$-\mathcal T_A r/2$. Project onto $h$, use
$\langle h,r_\perp\rangle_1=\langle h,k_A\rangle_1=0$, and add $\ell'$.
\end{proof}
If all $A$ inputs satisfy $|\theta|\le\theta_0<\pi/2$ and
$S_A=\int_A|U|^2\dd\lambda_0$, then
\begin{equation}\label{pa:trackingdamping}
 \kappa_A\ge\frac{Q_1}{2}\cos^2\theta_0
 \Phi(\cos\theta_0/q)^2S_A.
\end{equation}
Indeed
$\E_{P_1}[X_1(a^\top X)_+]=Q_1a_1\Phi(a_1/q)+q\phi(a_1/q)$,
and $\lambda_A=\int_A\langle h,(U^\top X)_+\rangle_1^2\dd\lambda_0$.
If $\kappa_A\ge\kappa_*>0$ and $|\eta_A|\le\eta_*$, variation of constants gives
\[
 |\delta(\tau)|\le e^{-\kappa_*(\tau-\tau_0)}|\delta(\tau_0)|
 +\frac{\eta_*}{\kappa_*}(1-e^{-\kappa_*(\tau-\tau_0)}).
\]
\openobligation{The useful forcing bound, the entry hypotheses, and a positive
signed lower bound on the lag are not supplied by this identity. Bounding
$|\eta_A|$ alone does not prove $\delta>0$. The fixed labels here are not the
instantaneous moving `main' mask in the diagnostics.}

\section{One noise-scaled population system}\label{pa:scaled}
\source{\cite{TN}, Sections 3--4, and \cite{SP}, Sections 1--3.}
\scope{The limiting equations are a derived leading system. Theorems below
about this system are exact within it. A uniform finite-noise field expansion
and an initialized trajectory transfer are still open.}

\subsection{Coordinates and Gaussian overlap}
Write $\lambda=\rho^2$, and define
\[
 \phi(r)=\frac{e^{-r^2/2}}{\sqrt{2\pi}},\quad
 \Phi(r)=\int_{-\infty}^r\phi(s)\dd s,\quad
 K(r)=\phi(r)+r\Phi(r),\quad g(r)=K(r)-r.
\]
For a standard normal $Z$, $K(r)=\E(Z+r)_+$,
$K'(r)=\Phi(r)$, and $g(r)=\E[-(Z+r)]_+>0$.
Use strong and weak scaled coordinates
\begin{align}\label{pa:charts}
 a_s&=(\cos(qx),\sin(qx)),&b_s&=(\cos(qy),\sin(qy)),\\
 a_w&=(\sin(qu),\cos(qu)),&b_w&=(-\sin(qv),\cos(qv)).\nonumber
\end{align}
A positive $u$ is a weak input tilt below $90^\circ$, toward the strong axis.
A positive $v$ is a weak output tilt above $90^\circ$.

\begin{lemma}[Overlap kernel]\label{pa:overlap}
For finite $t,r$,
\begin{equation}\label{pa:F}
 F(t,r):=\E[(Z+r)_+\1_{\{Z+t>0\}}]
 =\phi(\min(t,r))+r\Phi(\min(t,r)).
\end{equation}
Its first derivatives are
\[
 F_t(t,r)=(r-t)_+\phi(t),\qquad
 F_r(t,r)=\Phi(\min(t,r)).
\]
It is $C^{1,1}$ on compact sets but need not be $C^2$ at $t=r$.
\end{lemma}
\begin{proof}
The integration interval is $Z>-\min(t,r)$. Integrating $Z+r$ on this
interval gives \eqref{pa:F}. Differentiating its two branches gives the
stated derivatives, which agree on the diagonal. Their compact-set
Lipschitz property follows from the positive-part and minimum functions.
The derivative of $F_t$ changes branch at the diagonal, which prevents a
general classical second derivative there.
\end{proof}

Let $\mu_s(dx\,dy)$ and $\mu_w(du\,dv)$ be finite nonnegative squared-radius
measures, and use \emph{unnormalized} moments
\begin{align*}
 M&=\int\dd\mu_s,&N&=\int\dd\mu_w,&Y&=\int y\dd\mu_s,&V&=\int v\dd\mu_w,\\
 K_s&=\int K(\rho x)\dd\mu_s,&K_w&=\int K(u)\dd\mu_w.
\end{align*}
Define
\[
 \mathcal S(x)=\int F(\rho x,\rho\widetilde x)\dd\mu_s(\widetilde x,\widetilde y),
 \qquad
 \mathcal W(u)=\int F(u,\widetilde u)\dd\mu_w(\widetilde u,\widetilde v).
\]

\begin{definition}[Leading scaled population equations]\label{pa:leading}
The characteristic velocities are
\begin{align}
 x'={}&\tfrac12[-(1-M)x+\rho\phi(\rho x)-\rho\mathcal S(x)
 +\lambda\{V+(1-N)y\}\Phi(\rho x)],\label{pa:pop-x}\\
 y'={}&\tfrac12[-(1-M)y-Y-K_w+\rho(1-N)K(\rho x)],\label{pa:pop-y}\\
 u'={}&\tfrac12[\phi(u)-\mathcal W(u)-\{Y+(1-M)v\}\Phi(u)
 -\lambda(1-N)u],\label{pa:pop-u}\\
 v'={}&\tfrac12[\rho K_s-\lambda V-\lambda(1-N)v-(1-M)K(u)].\label{pa:pop-v}
\end{align}
The associated transport--reaction equations are
\begin{align}
 \partial_\tau\mu_s+\partial_x(x'\mu_s)+\partial_y(y'\mu_s)&=(1-M)\mu_s,
 \label{pa:poptransport}\\
 \partial_\tau\mu_w+\partial_u(u'\mu_w)+\partial_v(v'\mu_w)&=\lambda(1-N)\mu_w.
 \nonumber
\end{align}
We work with characteristic solutions on finite intervals with bounded
support and finite masses, or with explicitly sufficient moment bounds.
\end{definition}

\paragraph{Leading-order derivation, with the reduction gap exposed.}
On cluster 1 write $X=(1+qZ_1,qZ_2)$ and on cluster 2 write
$X=(qZ_1,\rho+qZ_2)$. For bounded scaled coordinates the own-cluster gates
are active except for Gaussian tails, while the cross gates remain
nontrivial. The leading residuals are
\begin{align*}
 P_1:\quad e_1(X)&=(1-M)(1+qZ_1)+O(q^2),\\
 e_2(X)&=q\left[Z_2-Y-\int(Z_2+u)_+\dd\mu_w\right]+O(q^2),\\
 P_2:\quad e_1(X)&=q\left[Z_1-\int(Z_1+\rho x)_+\dd\mu_s+\rho V\right]+O(q^2),\\
 e_2(X)&=(1-N)(\rho+qZ_2)+O(q^2).
\end{align*}
Insertion into the exact angular/radial field uses
$\E Z\1_{\{Z+r>0\}}=\phi(r)$ and Lemma~\ref{pa:overlap}.
It gives \eqref{pa:pop-x}--\eqref{pa:poptransport} as the leading coefficients.
These residual expansions describe the derivation; they do not alone prove
an $O(q^2)$ estimate for the \emph{scaled vector field}. In particular,
angular velocities are divided by $q$, gate indicators must be handled in
Gaussian coordinates, and long-time errors require a separate argument.
The tested $O(q^2)$ field accuracy is evidence, not a substitute for that lemma.
The isotropic population is not initially supported in \eqref{pa:charts};
uncharted labels and chart-entry flux cannot be omitted.

\begin{proposition}[Mass equations and atomic restrictions]\label{pa:masslaws}
For a closed pair of family measures satisfying
\eqref{pa:poptransport},
\begin{equation}\label{pa:logistic}
 M'=M(1-M),\qquad N'=\lambda N(1-N).
\end{equation}
The relative weights of fixed labels within either family are constant.
One atom in each family reduces the system to
\begin{align}
 m'&=m(1-m),&n'&=\lambda n(1-n),\label{pa:6d}\\
 x'&=\tfrac12[-(1-m)x+\rho(1-m)\phi(\rho x)
 +\lambda\{nv-mx+(1-n)y\}\Phi(\rho x)],\nonumber\\
 y'&=\tfrac12[-y-nK(u)+\rho(1-n)K(\rho x)],\nonumber\\
 u'&=\tfrac12[(1-n)\phi(u)-\{nu+my+(1-m)v\}\Phi(u)
 -\lambda(1-n)u],\nonumber\\
 v'&=\tfrac12[\rho mK(\rho x)-\lambda v-(1-m)K(u)].\nonumber
\end{align}
\end{proposition}
\begin{proof}
Integrate the transport equations. Every fixed strong label has the same
relative reaction rate $1-M$, and every fixed weak label has rate
$\lambda(1-N)$, proving the ratio assertion. For Dirac measures,
$\mathcal S(x)=mK(\rho x)$, $\mathcal W(u)=nK(u)$,
$Y=my$, and $V=nv$. Substitution and $K(r)=\phi(r)+r\Phi(r)$ give
\eqref{pa:6d}.
\end{proof}
The ratio assertion is not an exact finite-noise fact and does not hold for
a family redefined by a time-varying angular mask without flux terms.

The strong-family source decomposition is
\begin{align}
 x'_1&=-\tfrac12(1-M)x,&
 y'_1&=-\tfrac12[(1-M)y+Y+K_w],\label{pa:pop-sources}\\
 x'_2&=\tfrac12[\rho\phi(\rho x)-\rho\mathcal S(x)
 +\lambda\{V+(1-N)y\}\Phi(\rho x)],&
 y'_2&=\tfrac\rho2(1-N)K(\rho x).\nonumber
\end{align}
Strong mass growth is cluster-1-sourced and weak mass growth is
cluster-2-sourced to leading order. The one-half mixture factor is already
included.

\section{Resident dynamics, shape instability, and entry obstacles}
\label{pa:shape}
\source{\cite{SP}, Sections 4--9, including the distinction between coherent
and centered perturbations.}

\subsection{The invariant strong-learning curve}
\begin{lemma}[Resident direction]\label{pa:resident}
For every $0<\rho<1$ there is a unique $a=a_\rho>0$ satisfying
\begin{equation}\label{pa:root-a}
 a=\rho K(\rho a).
\end{equation}
With $N=0$, $\mu_s=M(\tau)\delta_{(a,a)}$ and $M'=M(1-M)$ form an
invariant curve of the leading population system.
\end{lemma}
\begin{proof}
The derivative of $a-\rho K(\rho a)$ is
$1-\rho^2\Phi(\rho a)\ge1-\rho^2>0$. The value at zero is negative and
it tends to infinity. On the proposed curve, substitution yields
$x'=(1-M)[-a+\rho K(\rho a)]/2=0$ and
$y'=[-a+\rho K(\rho a)]/2=0$.
\end{proof}
This is a reference solution, not a claim that the full initialized family
is already at that direction.
Put $P=\Phi(\rho a)$ and $\alpha=\lambda P/2$.

\begin{proposition}[Coherent and centered linearizations]\label{pa:spectra}
For normalized mean-zero strong-label offsets
$x_\beta=a+\varepsilon\xi_\beta$, $y_\beta=a+\varepsilon\eta_\beta$,
$\E\xi=\E\eta=0$, the first-order shape equation is
\begin{equation}\label{pa:Jsh-res}
 \frac{d}{d\tau}\begin{pmatrix}\xi\\\eta\end{pmatrix}
 =\begin{pmatrix}-(1-M)/2&\alpha\\\alpha&-(1-M)/2\end{pmatrix}
 \begin{pmatrix}\xi\\\eta\end{pmatrix}.
\end{equation}
The coherent mean perturbation instead has matrix
\begin{equation}\label{pa:Jmean}
 J_{\rm mean}(M)=
 \begin{pmatrix}-(1-M)/2-M\alpha&\alpha\\\alpha&-1/2\end{pmatrix}.
\end{equation}
At $M=1$ the coherent block is negative definite, while the shape rates
are $+\alpha$ and $-\alpha$.
\end{proposition}
\begin{proof}
At equal thresholds $F_t(t,t)=0$ and $F_r(t,t)=\Phi(t)$.
For a centered perturbation the first-order donor-field variation vanishes,
because it depends only on the mean offsets. Linearizing the receiving
variables gives \eqref{pa:Jsh-res}. For a coherent displacement, the donor
mean changes: differentiating the atomic equations in \eqref{pa:6d} at
$N=0$, $x=y=a$, yields \eqref{pa:Jmean}. At $M=1$ its trace is negative
and determinant $\alpha(1/2-\alpha)>0$, since $0<\alpha<1/2$.
The centered eigenvectors are $(1,1)$ and $(1,-1)$.
\end{proof}

\begin{remark}[Population multiplicity]
Every admissible centered profile in the $(1,1)$ direction is an unstable
first-order shape mode at the resident state. Thus ``weak growth and strong
splitting'' names two \emph{types} of instability, not a proved two-dimensional
unstable manifold of the measure flow. The coherent finite-dimensional
saddle spectrum cannot be used to infer stable condensation.
\end{remark}

\begin{proposition}[Moving shape matrix]\label{pa:moving-shape}
About a coherent state of masses $M,N$ and coordinates $(x,y,u,v)$, centered
strong perturbations have matrix
\begin{equation}\label{pa:Jmoving}
 J_{\rm sh}=
 \begin{pmatrix}-(1-M)/2+\beta&A_\mathrm{sh}\\
 A_\mathrm{sh}&-(1-M)/2\end{pmatrix},
\end{equation}
where
\[
 A_\mathrm{sh}=\frac\lambda2(1-N)\Phi(\rho x),\qquad
 \beta=\frac{\rho^3}{2}[Nv+(1-N)y-x]\phi(\rho x).
\]
\end{proposition}
\begin{proof}
The first-order donor variation vanishes for centered offsets, as above.
Differentiate the receiving coordinates in \eqref{pa:pop-x}--\eqref{pa:pop-y}.
At coincident strong inputs the receiver derivative of $\mathcal S$ is zero.
The derivative of $\rho\phi(\rho x)$ combines with that of the final
$\Phi(\rho x)$ factor to give $\beta$. The cross derivatives give
$A_\mathrm{sh}$ and the remaining diagonal term is $-(1-M)/2$.
\end{proof}
The frozen coefficient $\alpha(1-N)$ is a resident approximation.
The angle-dependent coefficient and $\beta$ must be retained during
weak learning; the eigendirections need not stay at $45^\circ$.

\subsection{Exact linear multipliers, not a nonlinear passage theorem}
\begin{proposition}[Linear resident passage factors]\label{pa:multipliers}
Along the resident curve, centered eigenmode amplitudes satisfy
\begin{align}
 \frac{h_\pm(\tau_1)}{h_\pm(\tau_0)}
 &=e^{\pm\alpha(\tau_1-\tau_0)}\sqrt{\frac{M_0}{M_1}},\label{pa:mult-time}\\
 \frac{h_+(M_1)}{h_+(M_0)}
 &=\left(\frac{M_1}{M_0}\right)^{\alpha-1/2}
 \left(\frac{1-M_0}{1-M_1}\right)^\alpha.\label{pa:mult-mass}
\end{align}
For the separately defined frozen resident weak-growth approximation
$h_+'=\alpha(1-N)h_+$, $N'=\lambda N(1-N)$,
\begin{equation}\label{pa:mult-weak}
 h_+(N)=h_+(N_0)(N/N_0)^{\alpha/\lambda}.
\end{equation}
\end{proposition}
\begin{proof}
Integrate the eigenvalue $-(1-M)/2\pm\alpha$. Logistic growth gives
$\int(1-M)\dd\tau=\log(M_1/M_0)$ and
$\tau_1-\tau_0=\log[M_1(1-M_0)/(M_0(1-M_1))]$.
For the last equation divide the two scalar differential equations.
\end{proof}
\openobligation{A nonlinear passage theorem must bound the moving matrix,
mean backreaction, weak angular relaxation, and population outside the chart.
The useful candidate shape parameter is
$\delta_{\rm ent}N_{\rm ent}^{-\alpha/\lambda}$, not just
$\delta_{\rm ent}$. If $N_{\rm ent}\asymp S_0^{c_0}$ and a strong-stage
shape is $O(S_0^{1/2-\alpha})$, the formal exponent is
$\eta=1/2-\alpha-c_0\alpha/\lambda$. The entry laws and resulting
phase threshold are not proved by the linear factors.}

\subsection{Ordering within a common population field}
\begin{proposition}[Same-field order preservation]\label{pa:order}
Assume $0\le N\le1$. Two strong characteristics in the same evolving
leading population field that are initially ordered in both coordinates
remain so on a chart-valid interval. A pairwise co-ordered strong population
has nonnegative normalized covariance $\Cov(x,y)$.
\end{proposition}
\begin{proof}
Holding the common field fixed,
\[
 \partial_y x'=\partial_x y'
 =\frac\lambda2(1-N)\Phi(\rho x)\ge0.
\]
At a first equality of one ordered coordinate, the other ordered coordinate
cannot make the difference point outward. Local Lipschitz regularity and
a strict-barrier approximation give the standard comparison argument.
For normalized label weights,
\[
 \Cov(x,y)=\frac12\E_{\beta,\gamma}
 [(x_\beta-x_\gamma)(y_\beta-y_\gamma)]\ge0.
\]
\end{proof}
This does not prove that isotropic labels have entered such a chart-ordered
family, or that two different populations are ordered relative to one
another. The bridge-induced field change is addressed in
Section~\ref{pa:bridge}.

\subsection{The weak attracting direction is created during learning}
\begin{proposition}[No initial weak-chart equilibrium]\label{pa:noweakeq}
For $0<\lambda<1$, the weak test-characteristic system at $M=N=0$ has no
finite equilibrium.
\end{proposition}
\begin{proof}
The $v$ equation requires $v=-K(u)/\lambda$. The remaining equation would be
\[
 \phi(u)+K(u)\Phi(u)/\lambda-\lambda u=0.
\]
For $u\le0$ its left side is positive. For $u>0$ it is decreasing in
$\lambda$, so is at least its value at $\lambda=1$, namely
$(1+\Phi(u))[K(u)-u]>0$.
\end{proof}
For frozen resident strong mass $M$, eliminating the weak output coordinate
instead gives
\begin{equation}\label{pa:frozen-weak}
 F_M(u)=\phi(u)-Ma\left(1+\frac{1-M}{\lambda}\right)\Phi(u)
 +\frac{(1-M)^2}{\lambda}K(u)\Phi(u)-\lambda u.
\end{equation}
The equations $F_M=\partial_uF_M=0$ specify a frozen-parameter saddle-node
candidate. Its reported numerical value near $M=0.4768$ at $\rho=2/3$ is
not a proved time-varying capture threshold. The weak seed must be obtained
from actual arrival/capture of initialized labels, not postulated as an
already formed atom.

\section{Population compensation lag and probe geometry}\label{pa:lagpop}
\source{\cite{SP}, Sections 10--12.}
\scope{The sign law in this section assumes the \emph{leading} surface $M=1$.
It is not an exact finite-$q$ positivity theorem. Neither family is assumed
to be a point mass.}

\begin{theorem}[Positive lag for a distributed weak family]\label{pa:lagtheorem}
Let a bounded-chart characteristic solution of the leading system have
$M=1$, $0\le N\le1$. Define
\[
 L_0=Y+K_w,
 \quad \eta_w=\mu_w/N\ (N>0),\quad Q(u)=\Phi(u),
\]
and interpret terms multiplied by $N$ continuously at $N=0$. Put
\begin{equation}\label{pa:Xi-def}
 \Xi=\E_{\eta_w}[Q\phi]+\E_{\eta_w}[K]\E_{\eta_w}[Q^2]
 -\E_{u,r\sim\eta_w}[Q(u)F(u,r)],
\end{equation}
where the two labels in the last expectation are independent. Then
\begin{align}\label{pa:lag-pop}
 L_0'={}&-\frac{1+N\E Q^2}{2}L_0
 +\frac\rho2(1-N)K_s\nonumber\\
 &+\frac{N(1-N)}2\{\lambda\E(K+\phi)+\E(Q\phi)\}
 +\frac{N^2}{2}\Xi.
\end{align}
For finite coordinates, $\Xi>0$. Hence a positive incoming lag stays positive
through every finite chart-valid interval.
\end{theorem}
\begin{proof}
The strong weights are stationary at $M=1$, so
\[
 Y'=\tfrac12[-Y-K_w+\rho(1-N)K_s].
\]
The normalized weak weights are also constant along labels. Differentiate
$K_w=N\E K(u)$, use $N'=\lambda N(1-N)$, and insert
\[
 u'=\tfrac12[\phi(u)-N\E_rF(u,r)-YQ(u)-\lambda(1-N)u].
\]
Substitute $Y=L_0-N\E K$. Group the terms in $L_0$, then use
$u\Phi(u)=K(u)-\phi(u)$; this gives \eqref{pa:lag-pop}.

For the sign, put $r_-=\min(u,r)$ and $r_+=\max(u,r)$.
Symmetrization of \eqref{pa:Xi-def} gives
\begin{align}\label{pa:Xi-positive}
 \Xi=\frac12\E_{u,r}\big[&g(r_+)
 \{\Phi(r_-)^2+\Phi(r_+)^2\}\nonumber\\
 &+[K(r_+)-K(r_-)]\Phi(r_+)[1-\Phi(r_+)]\big]>0.
\end{align}
To check this equality, take $u\le r$. Then
$F(u,r)=\phi(u)+r\Phi(u)$ and $F(r,u)=K(u)$.
Twice the symmetrized integrand reduces to
\[
 \Phi(r)\phi(r)+[K(r)-r]\Phi(u)^2
 -K(u)\Phi(r)[1-\Phi(r)].
\]
Adding and subtracting $K(r)\Phi(r)[1-\Phi(r)]$ produces
\eqref{pa:Xi-positive}. Since $g>0$ and $K$ is increasing, all forcing
terms in \eqref{pa:lag-pop} are nonnegative. The integrating-factor formula
preserves a positive initial value.
\end{proof}
For a weak Dirac measure the last forcing is
$\Xi=[K(u)-u]\Phi(u)^2$, recovering the two-atom identity.

\begin{proposition}[Population probe law]\label{pa:probe-pop}
Assume the entire state is represented by the two scaled-chart measures,
with compact support, bounded masses, and no outside population. Use $\xi_q=(\sin(q\zeta),-\cos(q\zeta))$ with bounded $\zeta$.
For sufficiently small positive $q$ the weak population is probe-inactive.
Set
\begin{equation}\label{pa:exposure}
 A_\zeta=\int(\zeta-x)_+\dd\mu_s,\qquad
 B_\zeta=\int y(\zeta-x)_+\dd\mu_s.
\end{equation}
The static expansion is
\begin{equation}\label{pa:pop-error}
 E_{\xi_q}-\frac12=q^2\mathcal E_\zeta+O(q^4),\qquad
 \mathcal E_\zeta=\frac12A_\zeta^2-\zeta A_\zeta+B_\zeta.
\end{equation}
Within the leading $M=1$ dynamics, the exact source rates are
\begin{align}
 d_1&=-\tfrac12A_\zeta L_0,\label{pa:pop-help}\\
 d_2&=\int_{x<\zeta}(\zeta-A_\zeta-y)x'_2\dd\mu_s
 +\frac\rho2(1-N)\int(\zeta-x)_+K(\rho x)\dd\mu_s,\label{pa:pop-harm}
\end{align}
and $\mathcal E_\zeta'=d_1+d_2$ wherever the chain rule applies.
In particular $L_0>0$ and $A_\zeta>0$ imply $d_1<0$.
\end{proposition}
\begin{proof}
For a weak input,
$a_w^\top\xi_q=-\cos(q(u+\zeta))<0$ on a bounded chart for small $q$.
For a strong input,
$a_s^\top\xi_q=\sin(q(\zeta-x))$. Uniform Taylor expansion, retaining the
positive part, gives
$f_1(\xi_q)=qA_\zeta+O(q^3)$ and
$f_2(\xi_q)=q^2B_\zeta+O(q^4)$. Substitution into
$E_\xi=1/2-\xi^\top f(\xi)+|f(\xi)|^2/2$ gives
\eqref{pa:pop-error}. At $M=1$ there is no leading strong mass reaction.
Equation~\eqref{pa:pop-sources} gives $x'_1=0$, $y'_1=-L_0/2$, proving
\eqref{pa:pop-help}. For source 2, differentiate $A_\zeta$ and $B_\zeta$:
\[
 A'_{\zeta,2}=-\int_{x<\zeta}x'_2\dd\mu_s,\quad
 B'_{\zeta,2}=\int[(\zeta-x)_+y'_2-y\1_{\{x<\zeta\}}x'_2]\dd\mu_s.
\]
This yields \eqref{pa:pop-harm}. Absolute continuity supplies the
almost-everywhere version when labels cross a probe boundary.
\end{proof}
If outside labels remain, their output and rate must be retained as separate
remainders. The $O(q^4)$ assertion is static on the stated bounded chart,
not a uniform approximation of the initialized trajectory over a
logarithmically long delay. The second term of $d_2$ is retained until weak saturation.
No arbitrary-state positivity of $d_2$ is asserted.

\begin{proposition}[Active covariance and inactive compensation]\label{pa:active-error}
Assume $M=1$. If all strong labels are active,
\[
 \mathcal E_\zeta=-\frac12\zeta^2+\frac12\bar x^2
 +\bar y(\zeta-\bar x)-\Cov(x,y).
\]
In general let $A=\{x<\zeta\}$ have mass $M_A>0$, normalized means
$\bar x_A,\bar y_A$, and normalized covariance $C_A$. Then, with
$h=\zeta-\bar x_A$,
\begin{equation}\label{pa:partial-error}
 \mathcal E_\zeta=\frac12M_A^2h^2
 -M_Ah(\zeta-\bar y_A)-M_AC_A.
\end{equation}
If $B=A^c$ in the strong population and $Y_B=\int_B y\dd\mu_s$, then
\begin{equation}\label{pa:bridge-error}
 \mathcal E_\zeta=\frac12A_\zeta^2-\zeta A_\zeta
 +h(L_0-K_w-Y_B)-M_AC_A.
\end{equation}
\end{proposition}
\begin{proof}
One has $A_\zeta=M_Ah$ and
$B_\zeta=M_A\bar y_Ah-M_AC_A$. Substitute into
\eqref{pa:pop-error}. For the last formula use
$M_A\bar y_A=Y-Y_B=L_0-K_w-Y_B$.
\end{proof}
Positive covariance improves error at fixed mass and barycentre. It does
not determine the sign of a trained-depth change, because the mean,
active mass, and weak residual also change. Equation~\eqref{pa:bridge-error}
shows how a probe-inactive bridge enters compensation without contributing
direct output at that probe.

\section{The selected coherent orbit and an analytical reversal theorem}
\label{pa:orbit}
\source{\cite{TN}, Sections 6--9, with the real-root and integrable-negative-part
corrections in \cite{SP}, Section 9.}
\scope{This section proves reversal for an analytically defined orbit of the
leading system. It does not prove that the isotropic population selects this
orbit. In particular, centered splitting is absent from this invariant
coherent restriction, not assumed stable in the full population.}

On the $m=1$ surface of \eqref{pa:6d}, put $n=N$ and retain weak growth:
\begin{align}
 n'&=\lambda n(1-n),\label{pa:5d}\\
 x'&=\tfrac\lambda2[nv-x+(1-n)y]\Phi(\rho x),\nonumber\\
 y'&=\tfrac12[-y-nK(u)+\rho(1-n)K(\rho x)],\nonumber\\
 u'&=\tfrac12[(1-n)(\phi(u)-\lambda u)-(nu+y)\Phi(u)],\nonumber\\
 v'&=\tfrac12[\rho K(\rho x)-\lambda v].\nonumber
\end{align}
Setting $n=1$ in advance would remove the weak-learning transition.

\begin{lemma}[Real weak resident coordinate]\label{pa:weakroot}
For every $0<\rho<1$, the equation
\begin{equation}\label{pa:weakroot-eq}
 \phi(u_*)-a_\rho\Phi(u_*)-\lambda u_*=0
\end{equation}
has exactly one real root. The derivative of its left-hand side at the root
is negative. The root need not be positive.
\end{lemma}
\begin{proof}
Divide by $\Phi(u)>0$ and put $r(u)=\phi(u)/\Phi(u)$. Then
\[
 r'(u)=-r(u)[u+r(u)]<0,
\]
because $u+r(u)=K(u)/\Phi(u)>0$. Also
\[
 \left(\frac u{\Phi(u)}\right)'
 =\frac{\Phi(u)-u\phi(u)}{\Phi(u)^2}>0.
\]
For $u<0$ its numerator is positive directly; for $u\ge0$ it equals
$1/2$ at zero and its derivative is $u^2\phi(u)\ge0$.
Thus $r(u)-a_\rho-\lambda u/\Phi(u)$ is strictly decreasing from
$+\infty$ to $-\infty$. At its unique zero the derivative of the original
left side is $\Phi(u_*)$ times this negative derivative.
\end{proof}

\begin{proposition}[Analytically specified weak-growth branch]\label{pa:branch}
The point
\[
 P_\rho=(n,x,y,u,v)=(0,a,a,u_*,a/\lambda)
\]
is a hyperbolic saddle of \eqref{pa:5d}, with one positive eigenvalue
$\lambda$ and four negative eigenvalues. Its positive-$n$ unstable branch
is unique up to time translation. Fix its translation by
\begin{equation}\label{pa:weakclock}
 n(s)=\frac1{1+e^{-\lambda s}},
\end{equation}
so that $s=0$ is weak half-mass. This specifies the selected coherent orbit
without fitted entry coordinates.
\end{proposition}
\begin{proof}
The resident equations follow from \eqref{pa:root-a} and
\eqref{pa:weakroot-eq}. The strong angular block is
$\bigl(\begin{smallmatrix}-\alpha&\alpha\\\alpha&-1/2\end{smallmatrix}\bigr)$,
negative definite by Proposition~\ref{pa:spectra}. The remaining angular
eigenvalues are
\[
 -\tfrac12[(u_*+a)\phi(u_*)+\lambda],\qquad -\lambda/2.
\]
The first is negative by Lemma~\ref{pa:weakroot}. The Jacobian is block
triangular when weak mass is ordered first, and its weak-mass eigenvalue is
$\lambda$. The finite-dimensional unstable-manifold theorem applies to
this smooth vector field; all its hyperbolicity hypotheses have just been
checked. The positive branch has logistic weak mass, and
\eqref{pa:weakclock} fixes its translation. The continued branch stays in
$0<n<1$. Its forward existence is proved below.
\end{proof}
No infinite-dimensional unstable-manifold result is invoked for the full
measure flow. That flow has the additional shape sector in
Proposition~\ref{pa:spectra}.

For an active strong atom the exact finite-$q$ error at a two-atom state is
\begin{equation}\label{pa:atomicexacterror}
 E_{\xi_q}-\tfrac12
 =\frac{m^2}{2}\sin^2(q(\zeta-x))
 -m\sin(q(\zeta-x))\sin(q(\zeta-y)),
\end{equation}
provided the weak atom is inactive. It follows simply by taking the norm
of the single active output. At $m=1$ its leading observable is
\begin{equation}\label{pa:atomicerror}
 \mathcal E_\zeta=-\frac12\zeta^2+\frac12x^2+y(\zeta-x),
 \qquad \mathcal E_\zeta'=(x-y)x'+(\zeta-x)y'.
\end{equation}
For all $x$ use
$\mathcal E_\zeta=\tfrac12(\zeta-x)_+^2-(\zeta-x)_+(\zeta-y)$.

The population lag theorem becomes, with $L=y+nK(u)$ and $Q=\Phi(u)$,
\begin{align}\label{pa:atomiclag}
 L'={}&-\tfrac12(1+nQ^2)L+\tfrac\rho2(1-n)K(\rho x)\nonumber\\
 &+\tfrac{n(1-n)}2[\lambda(K(u)+\phi(u))+Q\phi(u)]
 +\tfrac{n^2}{2}[K(u)-u]Q^2.
\end{align}
Since the resident value is $a>0$, the selected branch has $L>0$. Its
source rates in the active phase are
\begin{align}
 d_1&=-\tfrac12(\zeta-x)L,\label{pa:atomicrates}\\
 d_2&=\tfrac\lambda2(x-y)[nv-x+(1-n)y]\Phi(\rho x)
 +\tfrac\rho2(\zeta-x)(1-n)K(\rho x),\nonumber
\end{align}
with $\mathcal E_\zeta'=d_1+d_2$. In particular the second source contains
output rotation before weak saturation.

\begin{theorem}[Analytical reversal on the selected limit orbit]\label{pa:limitreversal}
Fix $0<\rho<1$ and $\zeta>a_\rho$. On the positive weak-growth branch of
Proposition~\ref{pa:branch}, there are finite earlier times with
$\mathcal E_\zeta'<0$ and later times with $\mathcal E_\zeta'>0$, before
the first hit $x=\zeta$. There is a sign-changing zero of the total rate
in the active phase, and a neighborhood of that zero on which
$d_1<0<d_2$. Shorter ordered intervals have strictly negative and strictly
positive total rates. The error has a positive quantitative rebound; an
interior late witness can be chosen with rebound at least
$(\zeta-a_\rho)^2/4$ from an earlier attained minimum.
\end{theorem}
\begin{proof}
\emph{1. Forward growth control.}
Let $R(s)=\max(|x|,|y|,|u|,|v|)$ after a finite initial time $s_0$ on the
branch. Use $0<K(r)\le r_++\phi(0)$. At an active maximum of $|x|$, the
convex combination $nv+(1-n)y$ and the negative self term imply an upper
Dini derivative at most zero. At a maximum of $|y|$ or $|v|$, the upper
bound is $\phi(0)/2$. At a maximum of $|u|$, the bound is
$(1-n)[\phi(0)+(1-\lambda)R]/2$. Therefore
\[
 D^+R\le\frac{\phi(0)}2+\frac{1-\lambda}{2}(1-n)R.
\]
The coefficient $1-n$ is integrable forward by \eqref{pa:weakclock}.
Gronwall gives $R(s)=O(1+s-s_0)$, hence no finite-time blowup and
$(1-n)R\to0$.

\emph{2. Strict initial improvement along the weak branch.}
Expand the branch at $n=0$:
\[
 x=a+X_1n+O(n^2),\qquad y=a+Y_1n+O(n^2).
\]
The invariance equation is
\[
 \begin{pmatrix}\lambda+\alpha&-\alpha\\-\alpha&\lambda+1/2\end{pmatrix}
 \begin{pmatrix}X_1\\Y_1\end{pmatrix}
 =\begin{pmatrix}a(1-\lambda)P/2\\-(a+K(u_*))/2\end{pmatrix}.
\]
For $\Delta=(\lambda+\alpha)(\lambda+1/2)-\alpha^2>0$,
\[
 Y_1=-\frac{\lambda}{4\Delta}
 [(2+P)(a+K(u_*))-a(1-\lambda)P^2]<0.
\]
The sign follows from $K(u_*)>0$ and $(1-\lambda)P^2<1<2+P$.
Since $x-y=O(n)$,
\[
 \mathcal E_\zeta'=(\zeta-a)\lambda Y_1n+O(n^2)<0
\]
for sufficiently small positive $n$. The resident limiting error is
$-(\zeta-a)^2/2$.

\emph{3. Invariants and a finite gate hit.}
The inequalities $v>0$, $v-y>0$, and $v-x>0$ hold at the resident point
and remain strict. At $v=0$, $v'=\rho K(\rho x)/2>0$.
At a possible boundary $v=y>0$,
\[
 (v-y)'=\tfrac12[\rho nK(\rho x)+(1-\lambda)y+nK(u)]>0.
\]
At $v=x$, the previous inequality implies $x-y>0$, and
\[
 (v-x)'=\tfrac\lambda2
 [K(\rho x)/\rho-x+(1-n)(x-y)\Phi(\rho x)]>0.
\]
Suppose for contradiction that $x<\zeta$ forever. Let $d=v-x>0$.
The negative part of $x'$ obeys
\[
 (x')_-\le\tfrac\lambda2(1-n)(v-y),
\]
which is integrable by step 1. Thus $x$ is bounded below, and under the
contradictory supposition it lies in a compact interval $[x_-,\zeta]$.
A direct subtraction gives
\[
 d'=\tfrac\lambda2\left[
 g_\rho(x)-(1+\Phi(\rho x))d
 +(1-n)(v-y)\Phi(\rho x)\right],
\]
where $g_\rho(x)=K(\rho x)/\rho-x>0$ is decreasing.
Consequently $d'\ge\lambda[g_\rho(\zeta)-2d]/2$, which gives an
eventual strictly positive lower bound on $d$. Since
$(1-n)(v-y)\to0$,
\[
 x'=\tfrac\lambda2[d-(1-n)(v-y)]\Phi(\rho x)
\]
eventually has a strictly positive lower bound. This contradicts $x<\zeta$.
There is therefore a finite first hit.

\emph{4. Reversal and source signs.}
At the first hit $\mathcal E_\zeta=0$. Before it, the error is negative
and decreases at some times, so it must increase at a later interior
point. The finite-dimensional vector field and active-phase observable
are real analytic in time. Their derivative is not identically zero;
a sign-changing zero can be chosen between a negative and a positive
rate. At that zero, $d_1<0$ by the positive-lag theorem, so
$d_2=-d_1>0$. Continuity gives the common sign neighborhood and two
strict total-rate subintervals.

Choose a finite early point with error below the resident limit. The
minimum before the hit is below $-(\zeta-a)^2/2$; continuity as the error
approaches zero at the hit permits a strictly interior late point with
rebound greater than $(\zeta-a)^2/4$. This proves the stated lower bound
without transferring the gate-closing point itself.
\end{proof}

\paragraph{Limits of this theorem.}
The weak mass is positive at every finite witness, but a specified numerical
learning threshold (for example $\mathsf m_2\ge1/2$ near the chosen crossover)
requires another quantitative argument and a mass-to-response comparison.
The theorem does not prove a unique or transverse crossing, the numerical
$1.1$ physical-time delay, common witness times across a large parameter
box, finite-noise error bounds, or canonical parameter coverage. Strict
witnesses at a fixed parameter can be extended to a local parameter/probe
neighborhood once the required continuous dependence and gate margins
are stated. Positivity of $n$ alone must not be presented as a proof that
weak learning is nearly complete.

\subsection{The additional fully learned reduction}
At $m=n=1$, put $c=-y$. Then
\[
 x'=\tfrac\lambda2(v-x)\Phi(\rho x),\quad
 v'=\tfrac\rho2[K(\rho x)-\rho v],\quad
 u'=\tfrac12(c-u)\Phi(u),\quad c'=\tfrac12[K(u)-c].
\]
After rescaling the first pair by $\rho$, both pairs have form
\[
 r'=\tfrac\gamma2(z-r)\Phi(r),\qquad
 z'=\tfrac\gamma2[K(r)-z].
\]
The region $0\le r<z<K(r)$ is forward invariant: at $z=r$ the gap
increases by $\gamma[K(r)-r]/2>0$, while at $z=K(r)$ the other gap
increases by $\gamma\Phi(r)^2(z-r)/2>0$. Within it, both coordinates
increase and cannot have a finite common equilibrium. This is a dynamical
Mills-gap explanation, not an instantaneous substitution $z=K(r)$.
The weak-learning equations must still retain $1-n$ terms before saturation.

\section{Exact finite-noise lag and non-small bridge refinements}
\label{pa:finiteq}
\source{The subsequent feedback and derivations \cite{C}, informed by the
observables implemented in \cite{D}. The formulas proved in this section
are mathematical identities or explicitly scoped comparisons, not consequences
of the numerical tests.}

\subsection{Three distinct compensation coordinates}
Retain $r_{21}$ from Proposition~\ref{pa:projection}. Define a different,
exact finite-noise mean lag by
\begin{equation}\label{pa:Lexdef}
 L_{1,q}^{\rm mean}:=\frac1q\E_{P_1}[f_2(X)-X_2]
 =\frac1q\E_{P_1}f_2(X).
\end{equation}
This is not $r_{21}/q$ and is not the leading quantity $L_0=Y+K_w$.

\begin{proposition}[Exact finite-noise mean lag]\label{pa:Lex}
For the full finite-mass angular state of the original population,
\begin{equation}\label{pa:Lexformula}
 L_{1,q}^{\rm mean}=\int b_2K(a_1/q)\dd\nu.
\end{equation}
Moreover,
\begin{equation}\label{pa:stein}
 \frac{1+q^2}{q}r_{21}
 =L_{1,q}^{\rm mean}+q\E_{P_1}[\partial_1 f_2(X)].
\end{equation}
If the state is partitioned into the scaled strong and weak charts plus a
remainder population, then exactly
\begin{align}\label{pa:Lexcharts}
 L_{1,q}^{\rm mean}={}&
 \int_s\sin(qy)K\!\left(\frac{\cos(qx)}q\right)\dd\mu_s\nonumber\\
 &+\int_w\cos(qv)K\!\left(\frac{\sin(qu)}q\right)\dd\mu_w
 +L_{\rm outside}.
\end{align}
\end{proposition}
\begin{proof}
For a unit direction $a$, $a^\top X\sim\mathcal N(a_1,q^2)$ under $P_1$,
so $\E(a^\top X)_+=qK(a_1/q)$. Insert the feature integral to prove
\eqref{pa:Lexformula}. Gaussian integration by parts gives
$\E[X_1f_2]=\E f_2+q^2\E\partial_1f_2$. The network is Lipschitz with
integrable weak derivative, so that identity is applicable. Since
$\E X_1X_2=0$, division by $q$ proves \eqref{pa:stein}. Substitution of
\eqref{pa:charts} proves \eqref{pa:Lexcharts}, without discarding the
remaining labels.
\end{proof}

\begin{proposition}[An explicit lag approximation bound]\label{pa:lagerror}
For the state partition of \eqref{pa:Lexcharts}, define
$L_0=\int_s y\dd\mu_s+\int_wK(u)\dd\mu_w$ at this same state.
Whenever the displayed moments are finite,
\begin{align}\label{pa:lagerror-eq}
 |L_{1,q}^{\rm mean}-L_0|\le{}&
 q^2\int_s\left(\frac{|y|^3}{6}+\frac{|y|x^2}{2}\right)\dd\mu_s\nonumber\\
 &+q^2\int_w\left(\frac{v^2K(u)}2+\frac{|u|^3}{6}\right)\dd\mu_w\nonumber\\
 &+\int_s|\sin(qy)|g\!\left(\frac{\cos(qx)}q\right)\dd\mu_s
 +|L_{\rm outside}|.
\end{align}
\end{proposition}
\begin{proof}
Use $K(t)=t+g(t)$ for the strong term. Then
\[
 \left|\frac{\sin(qy)\cos(qx)}q-y\right|
 \le q^2\left(|y|^3/6+|y|x^2/2\right).
\]
For the weak term, $|K(r)-K(u)|\le|r-u|$ because $0<\Phi<1$.
Writing $r=\sin(qu)/q$, adding and subtracting $\cos(qv)K(u)$ gives
\[
 |\cos(qv)K(r)-K(u)|
 \le q^2|u|^3/6+q^2v^2K(u)/2.
\]
These estimates use $|\sin z-z|\le|z|^3/6$,
$|1-\cos z|\le z^2/2$, and $|\cos z|\le1$. Integrating and keeping
the strong $g$ term and outside contribution proves the bound.
\end{proof}
The own-gate correction is small when the strong first coordinates are
uniformly positive, but the $q^2$ coefficients involve high angular
moments. A small mean or standard deviation alone does not control them.
Near a cancellation-dominated lag, an error small relative to the separate
positive and negative outputs may still be much larger than the lag.

\openobligation{The finite-noise helpful-rate target is
\[
 \frac{D_1^{\rm phys}}{\mu_1^2q^2}
 =-\frac12A_{q,\xi}L_{1,q}^{\rm mean}+\mathcal R_{1,q},
\]
with a precisely defined exposure, an analytically controlled remainder,
and a proved positive lag margin exceeding that remainder. The numerical
success of substituting the exact lag does not make the remainder zero.
A possible exact exposure convention is
$A_{q,\xi}=q^{-1}\int_{C_s}(a^\top\xi)_+\dd\nu$ for a specified strong
population $C_s$; all other labels and all omitted geometric factors must
then be included in $\mathcal R_{1,q}$. Its needed bound has not been proved.}

\subsection{Exact finite-noise same-field cross derivatives}
\begin{proposition}[Cross derivatives in the original angular field]
\label{pa:finiteqorder}
At a fixed continuous homogeneous residual field,
\[
 \partial_\psi\theta'=\partial_\theta\psi'
 =(b^\perp)^\top R(\theta)a^\perp.
\]
\end{proposition}
\begin{proof}
Differentiate \eqref{pa:exactflow}. The extra term in
$\partial_\theta\psi'$ is $(b^\perp)^\top R'(\theta)a$, which vanishes
because $R'(\theta)a=0$ by the Gaussian gate-boundary formula.
\end{proof}
Equality does not imply positivity. A finite-noise ordering or almost-ordering
argument must establish its sign or bound its negative contribution on a
region reached from initialization. This proposition fixes the field during
coordinate differentiation; it does not differentiate all other neurons'
responses to a state change.

\subsection{A nonperturbative instantaneous bridge-force lemma}
\label{pa:bridge}
\begin{proposition}[Upper-right donor displacement]\label{pa:bridge-force}
Consider the leading population field at a fixed state. Move strong mass
$\varepsilon>0$ from $(r_0,s_0)$ to $(r_B,s_B)$ with
$r_B\ge r_0$, $s_B\ge s_0$, keeping the total strong mass and the weak
measure fixed. Assume the resulting measure is nonnegative. At a fixed
receiving strong coordinate $(x,y)$,
\begin{align}
 \Delta x'&=-\frac{\varepsilon\rho}{2}
 [F(\rho x,\rho r_B)-F(\rho x,\rho r_0)]\le0,\label{pa:bridge-x}\\
 \Delta y'&=-\frac\varepsilon2(s_B-s_0)\le0.\label{pa:bridge-y}
\end{align}
If $x\le r_0$, then
\[
 \Delta x'=-\frac{\varepsilon\lambda}{2}(r_B-r_0)\Phi(\rho x).
\]
At a fixed receiving weak coordinate,
\begin{align}
 \Delta u'&=-\frac\varepsilon2(s_B-s_0)\Phi(u)\le0,\label{pa:bridge-u}\\
 \Delta v'&=\frac{\varepsilon\rho}{2}
 [K(\rho r_B)-K(\rho r_0)]\ge0.\label{pa:bridge-v}
\end{align}
\end{proposition}
\scope{Exact state-to-field comparison in the leading system. No
small-displacement assumption is made. It is not yet a comparison of two
self-consistently trained trajectories.}
\begin{proof}
In \eqref{pa:pop-x}, only $\mathcal S$ changes at the fixed receiver.
Lemma~\ref{pa:overlap} gives monotonicity in its donor argument, proving
\eqref{pa:bridge-x}; when $x\le r_0$ the two minimum arguments coincide
with $\rho x$, yielding the simpler expression. In \eqref{pa:pop-y},
only $Y$ changes, by $\varepsilon(s_B-s_0)$. In the two weak equations,
the changes are in $Y$ and $K_s$, respectively. Substitution proves the
remaining formulas.
\end{proof}
The first two signs implement the proposed exposure push and added output
compensation. The last two show why a trajectory comparison is coupled:
weak output compensation changes and feeds back through $V$ into the
strong input equation. It cannot be frozen for all later times without proof.

For distinct atomic strong labels $i,j$, the full field has donor derivatives
\begin{equation}\label{pa:donor-negative}
 \frac{\partial x_i'}{\partial x_j}
 =-\frac{\lambda m_j}{2}\Phi(\rho\min(x_i,x_j))<0,
 \qquad
 \frac{\partial y_i'}{\partial y_j}=-\frac{m_j}{2}<0.
\end{equation}
Thus same-field order preservation does not establish an order-preserving
flow in all particle coordinates. A main-block-versus-coherent-orbit
comparison needs differential inequalities for the changing common field.

\section{The small-splitting response: exact target functional, open sign}
\label{pa:response}
\source{\cite{SP}, Section 13. This is a small centered-shape calculation,
not a theorem for the canonical inactive bridge.}
\scope{The response equations and fixed-time coefficient are recorded with
their perturbative hypotheses. An analytical proof that the selected-orbit
coefficient is positive is open. Its numerical evaluation is not a proof.}

Let $M=1$, let the weak family be one atom, and let $(X,Y,U,V)$ be the
selected coherent orbit with prescribed logistic mass $N(s)$. For a bounded
profile $h_\gamma$ with $\E h=0$, $\E h^2=1$, take a centered strong
perturbation
\[
 x_\gamma=X+\varepsilon p(s)h_\gamma+O(\varepsilon^2),\qquad
 y_\gamma=Y+\varepsilon r(s)h_\gamma+O(\varepsilon^2).
\]
The first response obeys $(p,r)'=J_{\rm sh}(s)(p,r)$ with
\eqref{pa:Jmoving}. Write second-order means as
\[
 \bar x=X+\varepsilon^2X_2+o(\varepsilon^2),\quad
 \bar y=Y+\varepsilon^2Y_2+o(\varepsilon^2),
\]
and likewise $u=U+\varepsilon^2U_2+o(\varepsilon^2)$,
$v=V+\varepsilon^2V_2+o(\varepsilon^2)$. The perturbation family must specify
the initial second-order mean response; it is not an arbitrary free term.

Put $F_0=NV+(1-N)Y-X$, $P=\Phi(\rho X)$, $Q=\Phi(U)$. The coherent
angular Jacobian is
\begin{equation}\label{pa:Jresponse}
\setlength{\arraycolsep}{3pt}
 J_c=\begin{pmatrix}
 \frac\lambda2[-P+\rho F_0\phi(\rho X)]&\frac\lambda2(1-N)P&0&\frac\lambda2NP\\
 \frac\lambda2(1-N)P&-1/2&-NQ/2&0\\
 0&-Q/2&-\frac12[(U+Y)\phi(U)+NQ+\lambda(1-N)]&0\\
 \lambda P/2&0&0&-\lambda/2
 \end{pmatrix}.
\end{equation}

\begin{calculation}[Second-order mean response]\label{pa:response-eq}
On a bounded chart-valid finite interval, the second-order response obtained
by expanding the characteristic equations is
\begin{equation}\label{pa:response-system}
 \begin{pmatrix}X_2\\Y_2\\U_2\\V_2\end{pmatrix}'
 =J_c\begin{pmatrix}X_2\\Y_2\\U_2\\V_2\end{pmatrix}
 +\frac{\rho^3\phi(\rho X)}4
 \begin{pmatrix}
 -[1+\rho^2XF_0]p^2+2(1-N)pr\\
 (1-N)p^2\\0\\p^2
 \end{pmatrix}.
\end{equation}
\end{calculation}
\paragraph{Derivation and regularity guard.}
Do not use an unspecified classical Hessian of $F$ at coincident atoms.
For $\varepsilon>0$ and bounded offsets $a,b$, its one-sided expansion is
\[
 F(t+\varepsilon a,t+\varepsilon b)
 =K(t)+\varepsilon\Phi(t)b
 +\frac{\varepsilon^2\phi(t)}2[b^2-(b-a)_+^2]+O(\varepsilon^3).
\]
For independent identically distributed centered offsets,
$\E(b-a)_+^2=\Var(b)$. Hence the second-order double-averaged overlap
correction cancels at fixed mean. Taylor expanding the remaining smooth
terms of \eqref{pa:pop-x}--\eqref{pa:pop-v} gives the forcing in
\eqref{pa:response-system}. On finite bounded intervals, the explicit
one-sided remainder and locally Lipschitz first derivatives allow a
variation-of-constants justification. Passing from a finite entry time to
a selected incoming orbit at $s=-\infty$ additionally requires an incoming
response normalization and convergence estimate; numerical seeding at a
large negative time does not supply that estimate.

If all perturbed strong labels stay active and $s_*$ is an interior minimum
of the coherent observable, Proposition~\ref{pa:active-error} gives
\begin{equation}\label{pa:response-error}
 \mathcal E_\varepsilon(s_*)-\mathcal E_0(s_*)
 =\varepsilon^2[(X-Y)X_2+(\zeta-X)Y_2-pr]_{s_*}
 +o(\varepsilon^2).
\end{equation}
For input variance at weak half-mass
$\sigma_{\rm half}^2=\varepsilon^2p(0)^2$, define
\begin{equation}\label{pa:kappa}
 \kappa(\rho,\zeta)=
 \frac{[pr-(X-Y)X_2-(\zeta-X)Y_2]_{s_*}}{p(0)^2}.
\end{equation}
The direct covariance contribution is $pr$; the mean-response terms are of
the same order. Positivity does not follow solely from co-ordering or saddle
hyperbolicity. A negative correction evaluated at $s_*$ gives an improvement
lower bound at that point. An equality expansion of the \emph{new minimizing
value} additionally needs controlled minimizer selection, for example an
isolated nondegenerate minimum and a uniform expansion.

\openobligation{Prove the incoming response limit and $\kappa>0$ analytically
on a stated parameter region before advertising a positive enhancement
theorem. The numerical value $1.75053$ is a diagnostic evaluation. It is not
an exact universal constant, and the local expansion does not cover a
non-small inactive bridge. Basic reversal in Theorem A does not require
this optional enhancement result.}

\section{What still has to be proved for initialized A and B}\label{pa:roadmap}
\subsection{The theorem targets are still open}
\begin{target}[A: initialized two-concept signed competition]\label{pa:targetA}
For the exact flow initialized by \eqref{pa:init}, prove a nonempty explicitly
specified parameter region $\mathcal P$ with $\rho$ bounded away from zero,
a compact scaled-probe interval $I=[\zeta_-,\zeta_+]$ of positive length,
and a learned interval $J$ containing ordered positive-length subintervals
$J_-,J_+$, such that for all parameters and probes in the stated regime:
\begin{align*}
 \mathsf m_1,\mathsf m_2&\ge c_{\rm learn}>0,\qquad
 \mathcal L_p\le(1-\kappa_p)\mathcal L_p(0),\quad \kappa_p>0,\\
 D_1^\tau&\le-c_1q^2<0,\qquad D_2^\tau\ge c_2q^2>0\quad\text{on }J,\\
 D_1^\tau+D_2^\tau&\le-\gamma_-q^2\quad\text{on }J_-,\\
 D_1^\tau+D_2^\tau&\ge\gamma_+q^2\quad\text{on }J_+.
\end{align*}
The probe is $\xi_q=(\sin(q\zeta),-\cos(q\zeta))$.
The learning constants and rate margins must not vanish with $S_0$ or $q$
within the stated asymptotic regime. Prove a positive $cq^2$ rebound and
an angular sector of width at least $c_Iq$.
\end{target}
The cluster-loss reduction is a useful strengthened learning target; an
alternative response criterion must be justified as genuinely nontrivial.
The target does not require all-time opposite cluster signs, a unique
crossing, eventual recovery, permanent non-recovery, or an exact timing
constant. Those are separate obligations.

\paragraph{Canonical-facing versus asymptotic coverage.}
The scientific goal is a unified theory of the canonical parameters
$(\rho,q,S_0)=(2/3,0.05,2\cdot10^{-4})$ and the small-initialization
mechanism. An intermediate theorem for a rigorously specified small-noise
joint regime is valuable, but it must be labeled as intermediate unless
its estimates cover the canonical point. Do not replace the initialized
target with a conditional learned-state assumption or a dominant-concept
regime. The bridge and coherent edge belong to the same population model;
this fact alone is not a coverage proof.

\begin{target}[B: an initialized reciprocal-compensation explanation]
\label{pa:targetB}
Along the same initialized trajectory used for A, prove that Gaussian
cross-output leakage produces a compensating active response and that the
exact cluster-1 help admits a signed lag/exposure representation with
remainder smaller than its margin. Prove that the selected cluster-2
residual produces harmful input rotation of the probe-active population,
retaining other motion components. Account for the probe-inactive bridge
through its overlap and output moments, including the induced weak-family
feedback. Establish the temporal inequalities that make help dominate
first and harm dominate later. The identities and bounds in this proof
must imply A's crossover, rather than being fitted afterward to an
independently certified reversal.
\end{target}
The small-splitting coefficient is an optional refinement of B, not a
replacement for the large-bridge mechanism. A static covariance advantage
or donor-displacement sign is not yet the needed trajectory statement.

\subsection{Lemma roadmap with exact status}
The identifiers AN0--AN9 are roadmap identifiers, distinct from theorem
numbers in this document. Each new increment must state which identifier
it advances and which hypotheses remain unproved.

\paragraph{AN0: conventions and existing analytical tools (available).}
Use Section~\ref{pa:setup}, Propositions~\ref{pa:earlybounds}--\ref{pa:canonicalescape},
the exact identities in Section~\ref{pa:bookkeeping}, and the internally
proved leading-system results. The local scaled field is a derived model;
its finite-noise validity is not an initialized theorem. Read only the
relevant foundation propositions, not a new broad foundations program.

\paragraph{AN1: exact finite-noise lag and help closure (open).}
\emph{Input:} the exact flow and a quantitatively specified candidate
learning regime, without fitted trajectory data.
\emph{Required conclusion:} a dynamical identity or comparison inequality
for $L_{1,q}^{\rm mean}$ (or a precisely related exact observable), a positive
signed margin, and a helpful-rate remainder small relative to that margin.
\emph{Available:} Propositions~\ref{pa:projection}, \ref{pa:tracking},
\ref{pa:Lex}, and \ref{pa:lagerror}; the limiting positive-lag theorem.
\emph{Trap:} small coordinate loss bounds the absolute lag but not its sign;
$L_0$ may poorly approximate a small finite-$q$ cancellation.

\paragraph{AN2: global isotropic entry and weak capture (open; central gap).}
\emph{Input:} precisely \eqref{pa:init} and the exact Gaussian flow.
\emph{Required conclusion:} analytically defined phases/entry times,
substantial strong progress, a lower and upper bound for the weak incoming
seed, strong mean and shape bounds, weak angular bounds, and weighted
control of every population outside the selected charts. Derive ordering
or a quantified almost-order property when needed.
\emph{Available:} early escape, the target-only comparison, and the frozen
weak equilibrium equation. \emph{Trap:} the weak attractor is not present
in the initial scaled field; mass or angular windows measured from saved
states do not prove capture. Treat chart entry and moving membership
rather than assuming fixed two-family support from time zero.

\paragraph{AN3: nonlinear contraction and resident passage (open).}
\emph{Input:} AN2 entry data, including $N_{\rm ent}$ and the centered
strong profile. \emph{Required conclusion:} nonlinear bounds through weak
amplification that preserve a usable shape/residual class, with explicit
mean feedback, weak relaxation, and finite-noise corrections.
\emph{Available:} Propositions~\ref{pa:spectra}--\ref{pa:multipliers}.
\emph{Trap:} the exact linear multiplier is not a nonlinear estimate.
The conditional exponent
$1/2-\alpha-c_0\alpha/\lambda$ requires an actual entry law for $c_0$;
the numerical value near $0.072$ and proposed $\rho_c$ are not premises.

\paragraph{AN4: coupled non-small bridge comparison (open).}
\emph{Input:} a reached ordered or almost-ordered population class, with
active and inactive strong mass and weak residual moments.
\emph{Required conclusion:} evolve active exposure, bridge overlap/output
forcing, compensation, and weak feedback by coupled inequalities. Prove
that the regime survives for the needed time, rather than simply that a
bridge force has the desired sign at an isolated state.
\emph{Available:} Propositions~\ref{pa:active-error}, \ref{pa:order},
and \ref{pa:bridge-force}. \emph{Trap:} between-population comparison does
not follow from the same-field cooperative two-coordinate property.
The canonical bridge is not assumed to be a small variance perturbation.

\paragraph{AN5: substantial two-concept learning (open at the needed regime).}
\emph{Input:} the phase and residual estimates of AN2--AN4.
\emph{Required conclusion:} bounds on the exact observables
\eqref{pa:learning}, not merely mass labels. Give a positive response threshold
and a nontrivial loss reduction independent of vanishing initialization.
Do not infer monotonicity of $\mathsf m_p$ or $\mathcal L_p$ from total-loss
dissipation. A progress-clock hitting time requires existence and the
needed ordering/transversality proof.

\paragraph{AN6: harmful source sign and dominance switch (open for the
reached split population).}
\emph{Input:} exact lag/help bounds and the reached coupled population class.
\emph{Required conclusion:} a positive harmful rate, the complete remainder
budget, and a forced early-to-late crossover before first exit from the
regime. Preserve the distinction between training source, moving carrier,
and residual producer. \emph{Available:} the selected coherent reversal,
\eqref{pa:pop-harm}, \eqref{pa:producer}, and \eqref{pa:GLCidentity}.
The large-bridge statement is not proved by continuity from an atom without
a quantitative perturbation bound. Local opposite signs are sufficient;
the observed weak rate is not positive throughout training.

\paragraph{AN7: uniform finite-noise and long-time matching (open).}
Prove the needed scaled field and observable errors or work directly with
the exact finite-noise identities. If using the leading system, control
\[
 \frac{E_{q,S_0}-1/2}{q^2}-\mathcal E,
 \qquad \frac{D_{p,q,S_0}^\tau}{q^2}-d_p
\]
on the relevant centered intervals and probe sector, with errors below
analytic margins. Bound omitted labels and gate-strip contributions.
A possible joint scaling condition is $q^2\log(1/S_0)\to0$ if the actual
proved field errors support it; it is not by itself sufficient and is not
equivalent to $S_0\to0$ at fixed $q$. Preserve exact coincidence of labels
when estimating noise effects on splitting, instead of adding a spurious
noise-generated shape source.

\paragraph{AN8: initialized assembly, sector, and parameter region (open).}
Remove every intermediate entry assumption by composition of the prior
lemmas. Derive positive-width common time windows and a probe sector with
sufficient gate control, then integrate the exact rates. Specify the real
admissible parameter region and whether it includes the canonical point.
No sampled sign, numerical phase boundary, or unspecified ``small enough''
constant may be reported as canonical coverage.

\paragraph{AN9: optional refinements (open; not prerequisites for basic A).}
Analytical positivity of $\kappa$; broad-population enhancement inequalities;
a simple/unique crossing and a timing-error theorem; finite-horizon
persistence beyond the first rebound; finite-width/sample lifting;
ambient-dimension or many-concept extensions. None may be inserted into
the completed A statement until separately proved.

\subsection{A sufficient final assembly lemma}
\begin{lemma}[Analytical window-to-reversal implication]\label{pa:assembly}
Let $E_\xi$ be absolutely continuous and $E_\xi'=D_1^\tau+D_2^\tau$ almost
everywhere. Suppose a common interval contains
$J_-=[a,b]$ and $J_+=[c,d]$ with $a<b<c<d$, and, uniformly in the stated
probe/parameter set,
\[
 E_\xi'\le-\gamma_-q^2\ \text{a.e. on }[a,b],\qquad
 E_\xi'\ge\gamma_+q^2\ \text{a.e. on }[c,d],
 \quad\gamma_\pm>0.
\]
Then $E_\xi$ has a minimizer $\tau_*$ in $(a,d)$ and
\[
 E_\xi(a)-E_\xi(\tau_*)\ge\gamma_-q^2(b-a),\qquad
 E_\xi(d)-E_\xi(\tau_*)\ge\gamma_+q^2(d-c).
\]
If the individual signed-rate and learning assumptions of
Target~\ref{pa:targetA} have also been proved on the common interval, the
result has that exact attribution and learning content.
\end{lemma}
\begin{proof}
Integrate the derivative inequalities. Continuity on $[a,d]$ gives a
minimum. Since $E(b)<E(a)$ and $E(c)<E(d)$, a minimum cannot occur at
an endpoint. Comparing its value with $E(b)$ and $E(c)$ proves the two
bounds. Individual source signs and learning levels are separate premises;
the total-rate inequalities do not create them.
\end{proof}
This lemma is elementary and proved. Its hypotheses for the initialized
finite-noise population remain the substantive research task.

\appendix
\section{Empirical evidence: context only, never proof premises}\label{pa:evidence}
The following values are recorded to preserve the scientific target and to
avoid forgetting corrections between chats. They are numerical observations
or diagnostic evaluations from the supplied files and reports, not analytical
bounds on an initialized continuum trajectory.

\begin{longtable}{@{}p{0.37\textwidth}p{0.56\textwidth}@{}}
\toprule
Observation & Reported value and qualification\\
\midrule
\endhead
Canonical crossover & Fine-trace physical time about $3.675245$ at
$-85^\circ$; weak response about $0.99$. This does not prove a learning
threshold or exact timing theorem.\\
Early/late canonical rate signs & At physical time $3.5$ the total rate is
about $-7.93\cdot10^{-4}$; at $4$ it is about $+8.41\cdot10^{-4}$.
The weak rate is briefly negative near $2.75$, so opposite signs are a
learned-window conclusion, not an all-time claim.\\
Selected two-atom limit & Numerical scaled minimum about $-2.0092456$,
physical delay after weak half-mass about $1.10119$ for $\mu_1=3$.
The analytical result proves existence of reversal, not those decimals.\\
Split spectrum & At $\rho=2/3$, numerical evaluations give
$a\simeq0.351285$, $\alpha\simeq0.131685$, split instability threshold
$M\simeq0.736631$. The exact theorem uses the implicit root and
$M>1-2\alpha$.\\
Small-split response & Numerical $\kappa\simeq1.75053$, corroborated by
two-sub-atom calculations. Analytical positivity and a parameter interval
are open.\\
Conditional exponent & If the proposed seed exponent is $c_0=1$, the
formal $\eta$ is about $0.0720$ and its predicted zero near $\rho=0.7598$.
These are not proved population phase laws.\\
Finite-noise lag cancellation & The latest feedback reports canonical
$L^{\rm mean}_{1,q}\simeq0.059$ versus $L_0\simeq0.363$ at the crossing.
The exposure product fits far better with the exact lag. Positive and
negative lag contributions nearly cancel; their relative accuracy must not
be confused with lag-relative accuracy.\\
Canonical bridge enhancement & Latest reported canonical scaled error near
$-7.12$, bridge mass about $0.12$, and bridge output coordinate about $7.2$.
Direct active covariance accounts for a minority of excess depth in those
runs. This supports a non-small bridge comparison, not an extrapolation
of the small-split expansion.\\
Field and ordering checks & Multi-atom field tests show fourfold error
reduction under halving $q$ at selected states; sampled strong labels are
co-ordered. Neither is a uniform finite-noise or initialized-order theorem.\\
\bottomrule
\end{longtable}

\paragraph{Reporting guards.}
The new exact lag in \cite{D} includes all labels, while the leading moments
use angular masks. Its $u>3$ statistic refers to weak inputs below
$90^\circ$ by more than $3q$, and records absolute mass rather than a
normalized fraction. Density bins in old diagnostics are not pointwise
upper bounds. Median-centered RMS spread is not variance about the mean.
Counterfactual snapshots are not trained counterfactual trajectories, and
circular-mean collapse changes mean magnitudes as well as covariance.
A three-block collapse was found to improve one cluster-1 coordinate loss
while worsening the other; it does not isolate covariance alone.

The depth conventions also differ:
\[
 \frac{1/2-E_0}{q^2}=\frac{S_0/4-S_0^2/32}{q^2}.
\]
Thus improvement from initialization and depth below $1/2$ need not agree.
Nearby saved snapshots and fine-trace minima must not be identified as the
same time. Floating-point reporting precision is not a uniform sign proof.

\section{Review protocol and integration discipline}\label{pa:review}
Future work must be supplied as a standalone proof increment, for example
\path{AN02_isotropic_entry_v1.tex}, accompanied by
\path{AN02_isotropic_entry_v1_review.md}. The review note must state the
roadmap item, exact assumptions, conclusion, dependencies by label,
mathematical status, proof gaps or failed approaches, and proposed insertion
point in a future checkpoint. Give new labels a unique prefix such as
\verb|inc:AN02:v1:|. Do not overwrite a previous version or silently change
its hypotheses.

The current foundations and checkpoint stay read-only. Compilation and
formatting checks are allowed; they are not mathematical verification.
Numerical computation may be used only as explicitly labeled diagnostic
support, not to discharge any proof obligation. Do not introduce a
trajectory-certificate fallback. After explicit user approval of a checked
increment, create a new consolidated version with a changelog and an
old-to-new label map. Merely producing a proof draft does not authorize its
insertion into the master document.

\section{Source provenance}\label{pa:provenance}
The next table records source files used for this consolidation. Short hash
prefixes identify the exact local bytes; the companion handoff records full
SHA-256 values. Source files were not modified. The latest finite-noise-lag
and bridge-displacement derivations are taken from the conversation and
are consolidated here with explicit statements and proofs; their initialized
applications remain open.

\begin{longtable}{@{}p{0.66\textwidth}p{0.26\textwidth}@{}}
\toprule Source & SHA-256 prefix\\\midrule\endhead
\path{RELU_POPULATION_FOUNDATIONS_v2.tex} & \texttt{eefd557e0c725c87}\\
\path{RELU_POPULATION_FOUNDATIONS_v2.md} & \texttt{fed996786979ed10}\\
\path{LEAK_COMPENSATION_FACTORIZATION_AUDIT.md} & \texttt{6112f8fd549adab4}\\
\path{TWO_NEURON_SMALL_NOISE_LIMIT.md} & \texttt{7f4e2f74b83e0c8c}\\
\path{SCALED_POPULATION_SPLITTING_THEORY.md} & \texttt{68043ab7b75be2b6}\\
\path{theorem_A_initialized_escape.pdf} & \texttt{d7fedf43821b4cc1}\\
\path{SIM.md} & \texttt{672a8d492237c050}\\
\path{population_law_checks.py} & \texttt{32e6f576432c6e2f}\\
\path{split_mode_checks.py} & \texttt{411081436c9bccf9}\\
\bottomrule
\end{longtable}

\begingroup
\small
\begin{thebibliography}{99}
\bibitem[F]{F} \path{RELU_POPULATION_FOUNDATIONS_v2.tex}.
Existing exact population foundations, P1--P13 and supplements. Imported
mathematical base; not edited in this checkpoint.
\bibitem[E]{E} \path{theorem_A_initialized_escape.pdf}.
Analytical early-phase Propositions 1--3 and their proofs. Sections 4 onward
pursue the superseded certificate route and are excluded here.
\bibitem[LC]{LC} \path{LEAK_COMPENSATION_FACTORIZATION_AUDIT.md}.
Exact finite-parameter residual factorization and diagnostic audit.
Its numerical-continuation priorities are superseded by the analytical-only
roadmap; the exact identities are retained.
\bibitem[TN]{TN} \path{TWO_NEURON_SMALL_NOISE_LIMIT.md}.
Leading atomic equations and selected-orbit analytical reversal. The later
real-root and population-splitting corrections are incorporated explicitly.
\bibitem[SP]{SP} \path{SCALED_POPULATION_SPLITTING_THEORY.md}.
Leading population equation, resident shape spectrum, positive lag,
probe identities, and small-shape response calculation.
\bibitem[C]{C} Subsequent research conversation.
Fixed-label compensation tracking; exact finite-noise mean lag, Stein
relation and tail-moment bound; inactive-donor displacement and feedback;
clarification of within-population ordering versus between-population
comparison. These arguments are reproduced here, not assumed to establish
initialized A or B.
\bibitem[D]{D} \path{population_law_checks.py} and
\path{split_mode_checks.py}, with the user-supplied results.
Diagnostics and definitions of measured quantities only.
\bibitem[SIM]{SIM} \path{SIM.md} and the supplied predecessor paper,
\emph{Swing-by Dynamics in Concept Learning and Compositional Generalization}.
SIM task and historical growth--suppression setting. Its tied linear
matrix dynamics are not imported into the untied ReLU proof.
\end{thebibliography}
\endgroup
\end{document}
